% ============================================================
% 第12章 化学实验综合（必修实验基础 + 高三综合实验，对应广东卷第17题）
% ============================================================
% ---------- 本章 TikZ 样式与图元（前缀 c12） ----------
\tikzset{
  c12g/.style={line width=0.55pt},
  c12t/.style={line width=0.9pt,line cap=round},
  c12liq/.style={fill=mainblue!13},
  c12liqb/.style={fill=exorange!22},
  c12gas/.style={fill=green!45!yellow!22},
  c12stop/.style={line width=0.5pt,fill=black!28},
  c12flow/.style={-{Stealth[length=1.8mm,width=1.4mm]},line width=0.7pt,exorange},
  c12lab/.style={font=\scriptsize,align=center,inner sep=1pt},
  c12ptr/.style={line width=0.3pt,graytext},
  c12wire/.style={line width=0.8pt,exorange!70!black},
  % 集气瓶/洗气瓶：瓶底在原点，导管顶端在 y=2.4；左管长、右管短
  pics/c12bottle/.style={code={%
    \draw[c12g] (-0.25,1.58)--(-0.25,1.42)--(-0.45,1.25)--(-0.45,0.06) arc[start angle=180,end angle=270,radius=0.06]--(0.39,0) arc[start angle=270,end angle=360,radius=0.06]--(0.45,1.25)--(0.25,1.42)--(0.25,1.58);
    \draw[c12stop] (-0.29,1.5) rectangle (0.29,1.7);
    \draw[c12t] (-0.13,2.4)--(-0.13,0.15);
    \draw[c12t] (0.13,2.4)--(0.13,1.3);}},
  pics/c12wash/.style={code={\fill[c12liq] (-0.44,0.03) rectangle (0.44,1.0);\pic{c12bottle};}},
  pics/c12safe/.style={code={%
    \draw[c12g] (-0.25,1.58)--(-0.25,1.42)--(-0.45,1.25)--(-0.45,0.06) arc[start angle=180,end angle=270,radius=0.06]--(0.39,0) arc[start angle=270,end angle=360,radius=0.06]--(0.45,1.25)--(0.25,1.42)--(0.25,1.58);
    \draw[c12stop] (-0.29,1.5) rectangle (0.29,1.7);
    \draw[c12t] (-0.13,2.4)--(-0.13,1.3) (0.13,2.4)--(0.13,1.3);}},
  % 烧杯：杯底在原点
  pics/c12beaker/.style={code={%
    \fill[c12liq] (-0.49,0.02) rectangle (0.49,0.62);
    \draw[c12g] (-0.56,1.12)--(-0.5,1.06)--(-0.5,0)--(0.5,0)--(0.5,1.06)--(0.56,1.12);}},
  % 酒精灯：原点为被加热仪器底部
  pics/c12lamp/.style={code={%
    \draw[c12g,fill=black!5] (-0.42,-1.45)--(0.42,-1.45)--(0.42,-1.08) .. controls (0.42,-0.84) and (0.22,-0.78) .. (0.1,-0.76)--(-0.1,-0.76) .. controls (-0.22,-0.78) and (-0.42,-0.84) .. (-0.42,-1.08)--cycle;
    \draw[c12g,fill=black!25] (-0.09,-0.76) rectangle (0.09,-0.66);
    \draw[line width=0.6pt] (0,-0.66)--(0,-0.58);
    \fill[exorange!70] (0,-0.58) .. controls (0.14,-0.5) and (0.1,-0.3) .. (0,-0.13) .. controls (-0.1,-0.3) and (-0.14,-0.5) .. (0,-0.58);}},
  pics/c12net/.style={code={\draw[line width=1.1pt,black!50] (-0.7,-0.03)--(0.7,-0.03);}},
  % 圆底烧瓶：瓶底在原点，塞子顶端 y=2.08
  pics/c12rbfliq/.style={code={\begin{scope}\clip (0,0.72) circle (0.72);\fill[c12liq] (-1,-0.1) rectangle (1,0.45);\end{scope}}},
  pics/c12rbf/.style={code={%
    \draw[c12g] (-0.17,1.95)--(-0.17,1.42) arc[start angle=103.66,end angle=436.34,radius=0.72]--(0.17,1.95);
    \draw[c12stop] (-0.21,1.9) rectangle (0.21,2.08);}},
  % 分液漏斗主体：原点为活塞中心
  pics/c12sepliq/.style={code={\begin{scope}
    \clip (0.04,0.07) .. controls (0.15,0.22) and (0.55,0.55) .. (0.5,0.95) .. controls (0.46,1.3) and (0.15,1.4) .. (0.1,1.45) -- (-0.1,1.45) .. controls (-0.15,1.4) and (-0.46,1.3) .. (-0.5,0.95) .. controls (-0.55,0.55) and (-0.15,0.22) .. (-0.04,0.07) -- cycle;
    \fill[c12liq] (-0.6,0) rectangle (0.6,0.85);\end{scope}}},
  pics/c12sepbody/.style={code={%
    \draw[c12g] (0.04,0.07) .. controls (0.15,0.22) and (0.55,0.55) .. (0.5,0.95) .. controls (0.46,1.3) and (0.15,1.4) .. (0.1,1.45) -- (0.1,1.6);
    \draw[c12g] (-0.04,0.07) .. controls (-0.15,0.22) and (-0.55,0.55) .. (-0.5,0.95) .. controls (-0.46,1.3) and (-0.15,1.4) .. (-0.1,1.45) -- (-0.1,1.6);
    \draw[c12stop] (-0.13,1.55) rectangle (0.13,1.72);
    \draw[c12g,fill=white] (-0.2,-0.045) rectangle (0.2,0.045);
    \draw[c12g,fill=white] (0.22,0) circle (0.055);}},
  % 锥形瓶：瓶底在原点，瓶口 y=1.3
  pics/c12cone/.style={code={%
    \begin{scope}\clip (-0.15,1.3)--(-0.15,0.95)--(-0.62,0)--(0.62,0)--(0.15,0.95)--(0.15,1.3)--cycle;
    \fill[c12coneliq] (-1,0) rectangle (1,0.38);\end{scope}
    \draw[c12g] (-0.15,1.3)--(-0.15,0.95)--(-0.62,0)--(0.62,0)--(0.15,0.95)--(0.15,1.3);}},
  c12coneliq/.style={fill=mainblue!13},
}

\chapter{化学实验综合}

本章把必修中的实验基础和高三的综合实验放在一起学习：先认识仪器、掌握基本操作，再学物质的分离提纯和检验，然后学习气体制备装置的连接、滴定和量气等定量实验及误差分析，最后学习控制变量、提出假设并设计方案的探究实验。广东卷选择题中每年有 2～3 道实验题（2024年第6、8、11题；2025年第6、11、15题；2026年第9、12题），第17题是 14 分的综合实验探究题，三年都由基础操作、定量滴定和探究设计组成：2024年探究沉淀转化的速率与平衡的建立，2025年测定苯甲酸的浓度与 $K_\text{a}$ 并设计测定 HClO 的 $K_\text{a}$，2026年探究铜氨配合物与酸的反应。做这类题，除了会操作、会计算，还要能用规范的语言写出操作、现象和结论。

\section{常用仪器与基本操作}

\subsection{常用仪器}

认识一件仪器，抓住三个问题：能不能加热，量得多准，用来做什么。试管、蒸发皿、坩埚可以直接加热（灼烧固体用坩埚，放在泥三角上，用坩埚钳夹取）；烧杯、烧瓶、锥形瓶应垫陶土网或采用合适的加热浴，以免受热不均；量筒、容量瓶、滴定管等计量仪器不能用来加热。读数要依据仪器的分度值：托盘天平通常读到 0.1 g，10 mL 量筒通常读到 0.1 mL，滴定管读到 0.01 mL；量程较大的量筒分度更粗，不能把所有量筒都当作能读到 0.1 mL。量筒没有“0”刻度。图~\ref{fig:12-yiqi} 画出了常用玻璃仪器。

\tu{\begin{tikzpicture}[scale=0.92]
  % (a) 容量瓶
  \begin{scope}
    \clip (-0.12,2.3)--(-0.12,1.5) .. controls (-0.12,1.25) and (-0.62,1.05) .. (-0.62,0.55) .. controls (-0.62,0.18) .. (-0.45,0)--(0.45,0) .. controls (0.62,0.18) .. (0.62,0.55) .. controls (0.62,1.05) and (0.12,1.25) .. (0.12,1.5)--(0.12,2.3)--cycle;
    \fill[c12liq] (-1,0) rectangle (1,2.3);
  \end{scope}
  \draw[c12g] (-0.12,3.0)--(-0.12,1.5) .. controls (-0.12,1.25) and (-0.62,1.05) .. (-0.62,0.55) .. controls (-0.62,0.18) .. (-0.45,0)--(0.45,0) .. controls (0.62,0.18) .. (0.62,0.55) .. controls (0.62,1.05) and (0.12,1.25) .. (0.12,1.5)--(0.12,3.0);
  \draw[c12stop] (-0.16,2.95) rectangle (0.16,3.2);
  \draw[line width=0.7pt] (-0.12,2.3)--(0.12,2.3);
  \draw[c12ptr] (0.12,2.3)--(0.5,2.55) node[right,c12lab]{刻度线};
  \node[font=\tiny,align=center] at (0,0.55) {20\,℃\\250\,mL};
  \node[c12lab] at (0,-0.35) {(a) 容量瓶};
  % (b) 酸式滴定管
  \begin{scope}[xshift=2.1cm]
    \fill[c12liq] (-0.09,0.55) rectangle (0.09,3.0);
    \draw[c12g] (-0.09,3.7)--(-0.09,0.55) (0.09,3.7)--(0.09,0.55);
    \foreach \i in {0,...,6} \draw[line width=0.3pt] (-0.09,3.4-0.42*\i)--(0.0,3.4-0.42*\i);
    \node[font=\tiny,left] at (-0.09,3.4) {0};
    \draw[c12g,fill=white] (0,0.45) circle (0.1);
    \draw[line width=1pt] (0.1,0.45)--(0.32,0.45);
    \draw[c12g] (-0.04,0.36)--(-0.015,0.02) (0.04,0.36)--(0.015,0.02);
    \draw[c12ptr] (0.3,0.5)--(0.55,0.95) node[above,c12lab]{玻璃活塞};
    \node[c12lab] at (0,-0.35) {(b) 酸式滴定管};
  \end{scope}
  % (c) 碱式滴定管
  \begin{scope}[xshift=4.2cm]
    \fill[c12liq] (-0.09,0.78) rectangle (0.09,3.0);
    \draw[c12g] (-0.09,3.7)--(-0.09,0.78) (0.09,3.7)--(0.09,0.78);
    \foreach \i in {0,...,6} \draw[line width=0.3pt] (-0.09,3.4-0.42*\i)--(0.0,3.4-0.42*\i);
    \node[font=\tiny,left] at (-0.09,3.4) {0};
    \fill[black!55] (-0.075,0.33) rectangle (0.075,0.8);
    \draw[fill=white,line width=0.4pt] (0,0.56) circle (0.055);
    \draw[c12g] (-0.04,0.33)--(-0.015,0.02) (0.04,0.33)--(0.015,0.02);
    \draw[c12ptr] (0.08,0.56)--(0.42,0.95) node[above,c12lab]{橡胶管\\内有玻璃珠};
    \node[c12lab] at (0,-0.35) {(c) 碱式滴定管};
  \end{scope}
  % (d) 分液漏斗
  \begin{scope}[xshift=6.25cm]
    \draw[c12t] (0,0.05)--(0,1.15);
    \pic[scale=1.15] at (0,1.2) {c12sepliq};
    \pic[scale=1.15] at (0,1.2) {c12sepbody};
    \node[c12lab] at (0,-0.35) {(d) 分液漏斗};
  \end{scope}
  % (e) 恒压滴液漏斗
  \begin{scope}[xshift=8.35cm]
    \fill[c12liq] (-0.25,1.3) rectangle (0.25,2.2);
    \fill[c12liq] (-0.25,1.3)--(-0.04,1.05)--(0.04,1.05)--(0.25,1.3)--cycle;
    \draw[c12g] (-0.1,3.1)--(-0.1,2.9)--(-0.25,2.9)--(-0.25,1.3)--(-0.04,1.05)--(-0.04,0.98);
    \draw[c12g] (0.1,3.1)--(0.1,2.9)--(0.25,2.9)--(0.25,1.3)--(0.04,1.05)--(0.04,0.98);
    \draw[c12stop] (-0.13,3.05) rectangle (0.13,3.22);
    \draw[c12g,fill=white] (-0.2,0.89) rectangle (0.2,0.97);
    \draw[c12t] (0,0.88)--(0,0.05);
    \draw[c12t,mainblue!70!black] (0.25,2.75)--(0.5,2.75)--(0.5,0.5)--(0.02,0.5);
    \draw[c12ptr] (0.5,1.7)--(0.75,1.95) node[above,c12lab]{侧管};
    \node[c12lab] at (0,-0.35) {(e) 恒压滴液漏斗};
  \end{scope}
  % (f) 直形冷凝管
  \begin{scope}[xshift=10.6cm]
    \fill[c12liq] (-0.2,0.6) rectangle (0.2,3.2);
    \draw[c12g] (-0.2,0.6) rectangle (0.2,3.2);
    \draw[c12t,fill=white] (0,0.15)--(0,3.65);
    \draw[c12g] (0.2,0.85)--(0.45,0.85) (-0.2,2.95)--(-0.45,2.95);
    \draw[c12flow] (0.75,0.85)--(0.46,0.85);
    \node[c12lab,below] at (0.6,0.78) {进水};
    \draw[c12flow] (-0.46,2.95)--(-0.75,2.95);
    \node[c12lab,above] at (-0.6,3.02) {出水};
    \node[c12lab] at (0,-0.35) {(f) 直形冷凝管};
  \end{scope}
  % (g) 球形冷凝管
  \begin{scope}[xshift=12.6cm]
    \fill[c12liq] (-0.27,0.6) rectangle (0.27,3.2);
    \draw[c12g] (-0.27,0.6) rectangle (0.27,3.2);
    \draw[c12t] (0,0.15)--(0,0.83) (0,3.17)--(0,3.65);
    \foreach \y in {1.0,1.45,1.9,2.35,2.8} \draw[c12g,fill=white] (0,\y) circle (0.17);
    \foreach \y in {1.17,1.62,2.07,2.52} \draw[c12t] (0,\y)--(0,\y+0.11);
    \draw[c12g] (0.27,0.85)--(0.5,0.85) (-0.27,2.95)--(-0.5,2.95);
    \draw[c12flow] (0.8,0.85)--(0.51,0.85);
    \node[c12lab,below] at (0.65,0.78) {进水};
    \draw[c12flow] (-0.51,2.95)--(-0.8,2.95);
    \node[c12lab,above] at (-0.65,3.02) {出水};
    \node[c12lab] at (0,-0.35) {(g) 球形冷凝管};
  \end{scope}
\end{tikzpicture}}{几种常用玻璃仪器（简笔示意）\label{fig:12-yiqi}}

读图时注意每件仪器的结构特点，用途和使用要点都可以从结构推出来。

容量瓶（图~\ref{fig:12-yiqi}(a)）瓶颈上只有一条刻度线，瓶身标有温度和容积，所以它只能在标定温度下配制这一体积的溶液。使用前要检漏：加水、塞紧、倒立，看是否漏水；正立后把瓶塞旋转 180° 再倒立检查一次。容量瓶不能加热，也不能在瓶中溶解固体、稀释浓溶液或长期贮存溶液。溶解和稀释常伴有热效应，温度变了，溶液体积就和刻度不对应了。

滴定管（图(b)(c)）的“0”刻度在上方，从上往下读数增大，可读到 0.01 mL。酸式滴定管下端是玻璃活塞，可以盛多数酸性溶液（不能盛氢氟酸）和氧化性溶液（如 \ce{KMnO4}、\ce{I2} 溶液），不能盛碱性溶液，因为碱与玻璃中的 \ce{SiO2} 反应生成有黏性的硅酸盐，会把活塞粘住。碱式滴定管下端是内有玻璃珠的橡胶管，用手指挤压玻璃珠旁边的橡胶管，溶液就从缝隙中流出；它不能盛氧化性溶液，因为强氧化剂会腐蚀橡胶。现在实验室常用聚四氟乙烯活塞的滴定管，酸、碱溶液都能盛放。

分液漏斗（图(d)）用于萃取分液，也可以向反应容器中逐滴加入液体，使用前同样要检查活塞和上口玻璃塞是否漏水。恒压滴液漏斗（图(e)）多了一根侧管，把漏斗上方的空间和下方反应容器连通，两处气压相等，液体就能在重力作用下顺利流下。在密闭体系中滴加液体，用它还有一个好处：滴下的液体不会把等体积的气体“挤”出去，这一点在量气实验中很重要（见第~\ref{sec:12-dingliang}~节）。

冷凝管有两种。直形冷凝管（图(f)）用于蒸馏，冷凝液顺着向下倾斜的内管流入接收器；球形冷凝管（图(g)）的球泡增大冷却面积，竖直装在反应容器上方，使蒸气冷凝后流回容器，这叫冷凝回流。冷却水都应下口进、上口出，先保证夹套充满水。蒸馏时水与蒸气逆向流动，有利于充分换热；回流时蒸气总体向上、冷却水也向上，不能把“逆向”当作回流装置的理由。装置上端通常与大气相通，不能把受热体系完全封闭。

\subsection{基本操作}

下面几项操作在实验题中反复出现，记住步骤的同时要理解每一步的理由。

\begin{enumerate}
\item 检查气密性。原理是让装置内的气体压强发生变化，再看有没有相应的现象。微热法：把导管末端插入水中，用手捂热（或微热）容器，导管口有气泡冒出，停止加热后导管内形成一段稳定的水柱，说明不漏气。液差法：关闭止水夹，从长颈漏斗向容器中加水，使漏斗中的液面高于容器内的液面，一段时间后液面差不变，说明不漏气。
\item 测溶液的 pH。用洁净干燥的玻璃棒蘸取待测液点在 pH 试纸中部，与标准比色卡对照。试纸不能润湿，否则相当于把溶液稀释了；氯水、NaClO 溶液有漂白性，会把试纸漂白，不能用 pH 试纸测，要用 pH 计。
\item 洗涤沉淀并检验是否洗净。向过滤器中加蒸馏水至浸没沉淀，待水自然流下，重复 2～3 次。检验时取最后一次洗涤液少许于试管中，加入能与可溶性杂质离子反应的试剂，若无明显现象，说明已洗净。例如沉淀表面吸附 \ce{SO4^2-}，就滴加 \ce{BaCl2} 溶液检验。
\item 配制一定物质的量浓度的溶液。步骤为：计算、称量（或量取）、溶解（或稀释）并冷却至室温、转移、洗涤烧杯和玻璃棒 2～3 次并把洗涤液转入容量瓶、定容（加水至刻度线下 1～2 cm 时改用胶头滴管）、摇匀。分析误差只看两个量：由 $c=\dfrac{n}{V}$，凡使溶质 $n$ 偏小的操作使浓度偏低，凡使溶液体积 $V$ 偏小的操作使浓度偏高。
\item 安全措施。稀释浓硫酸时，把酸沿器壁缓慢注入水中并搅拌；加热液体前放入碎瓷片或沸石防暴沸；有毒气体实验在通风橱中进行，并设置尾气吸收。吸收 \ce{NH3}、HCl 等极易溶于水的气体时要防倒吸：倒扣漏斗的边缘只接触液面，不能深浸；也可用合适容积的安全瓶。不能把“装了漏斗”当作不再需要检查液面位置。
\end{enumerate}

\begin{liti}[基础]
用 NaOH 固体配制 250 mL 0.100\molL{} NaOH 溶液。(1) 用托盘天平称取 NaOH 的质量是多少？称量时 NaOH 应放在什么容器中？(2) 除托盘天平、药匙外，还需要哪些仪器？(3) 判断下列操作对所配溶液浓度的影响：①NaOH 溶解后未冷却就转移并定容；②定容时仰视刻度线；③摇匀后发现液面低于刻度线，又加水至刻度线；④容量瓶洗净后瓶内残留少量蒸馏水。
\tcblower
\jie (1) $m=0.250\times0.100\times40=1.0$ g。NaOH 易潮解且有腐蚀性，应放在小烧杯中称量，不能放在纸上。\\
(2) 烧杯、玻璃棒、250 mL 容量瓶、胶头滴管。\\
(3) ①NaOH 溶解放热，趁热定容，冷却后溶液体积缩小，$V$ 偏小，浓度偏高。②仰视时视线从下往上看刻度线，液面已经超过刻度线，$V$ 偏大，浓度偏低。③摇匀后部分溶液沾在瓶颈和瓶塞上，液面略低是正常的，再加水使 $V$ 偏大，浓度偏低。④定容时本来就要加水，残留的水不影响 $n$ 和 $V$，无影响。
NaOH 易吸水并吸收 \ce{CO2}，直接称量配得的是近似浓度；精确滴定所用 NaOH 标准溶液还需标定。
\end{liti}

判断一个装置能否达到实验目的，要依次检查四点：反应原理是否正确，试剂是否合适，装置是否合理（是否需要加热，导管长短，是否防倒吸），操作顺序是否正确。下面这道真题就是这样逐项检查的。

\begin{zhenti}{2024年广东卷第8题}
1810 年，化学家戴维首次确认“氯气”是一种新元素组成的单质。兴趣小组利用图~\ref{fig:12-z24-8} 所示装置进行实验，其中难以达到预期目的的是\xz
\tu{\begin{tikzpicture}[scale=0.8]
  % A
  \fill[black!70] (-0.3,0.12) circle (0.05) (-0.12,0.08) circle (0.05) (0.08,0.1) circle (0.05) (0.25,0.15) circle (0.05) (0.0,0.2) circle (0.05);
  \pic at (0,0) {c12rbf};
  \draw[c12t] (-0.06,0.9)--(-0.06,2.15);
  \pic at (-0.06,2.2) {c12sepliq};
  \pic at (-0.06,2.2) {c12sepbody};
  \draw[c12t] (0.09,1.6)--(0.09,2.0)--(1.0,2.0);
  \draw[c12flow] (0.55,2.15)--(1.0,2.15);
  \node[c12lab,left] at (-0.5,3.0) {盐酸};
  \draw[c12ptr] (-0.52,3.0)--(-0.3,2.8);
  \node[c12lab,left] at (-0.75,0.4) {\ce{MnO2}};
  \draw[c12ptr] (-0.75,0.4)--(-0.25,0.15);
  \node[c12lab] at (0.1,-0.4) {A．制备 \ce{Cl2}};
  % B
  \begin{scope}[xshift=3.0cm]
    \pic at (0,0) {c12wash};
    \draw[c12t] (-0.8,2.4)--(-0.13,2.4);
    \draw[c12flow] (-0.8,2.55)--(-0.35,2.55);
    \draw[c12t] (0.13,2.4)--(0.6,2.4);
    \draw[c12g] (0.6,2.46)--(0.75,2.46) (0.6,2.34)--(0.75,2.34);
    \draw[c12g] (0.75,2.46) arc[start angle=160,end angle=20,x radius=0.4,y radius=0.26] -- (1.5,2.46);
    \draw[c12g] (0.75,2.34) arc[start angle=-160,end angle=-20,x radius=0.4,y radius=0.26] -- (1.5,2.34);
    \foreach \x/\y in {0.92/2.42,1.05/2.5,1.12/2.33,1.25/2.45,1.0/2.3,1.32/2.36} \draw[line width=0.3pt] (\x,\y) circle (0.045);
    \node[c12lab] at (1.13,2.95) {无水 \ce{CaCl2}};
    \node[c12lab] at (0,-0.15) {饱和食盐水};
    \node[c12lab] at (0.4,-0.6) {B．净化、干燥 \ce{Cl2}};
  \end{scope}
  % C
  \begin{scope}[xshift=6.1cm]
    \pic at (0,0) {c12bottle};
    \draw[c12t] (-0.7,2.4)--(-0.13,2.4) (0.13,2.4)--(0.6,2.4);
    \draw[c12flow] (-0.7,2.55)--(-0.3,2.55);
    \node[c12lab,above] at (-0.55,2.6) {\ce{Cl2}};
    \draw[c12flow] (0.25,2.55)--(0.65,2.55);
    \node[c12lab] at (0,-0.6) {C．收集 \ce{Cl2}};
  \end{scope}
  % D
  \begin{scope}[xshift=8.4cm]
    \fill[c12gas] (-0.44,0.03) rectangle (0.44,1.25);
    \fill[c12gas] (-0.44,1.25)--(-0.25,1.42)--(0.25,1.42)--(0.44,1.25)--cycle;
    \draw[c12g] (-0.25,1.58)--(-0.25,1.42)--(-0.45,1.25)--(-0.45,0.06) arc[start angle=180,end angle=270,radius=0.06]--(0.39,0) arc[start angle=270,end angle=360,radius=0.06]--(0.45,1.25)--(0.25,1.42)--(0.25,1.58);
    \draw[c12t] (0,2.5)--(0,0.75);
    \fill[exorange!70] (0,0.72) .. controls (0.13,0.62) and (0.1,0.48) .. (0,0.36) .. controls (-0.1,0.48) and (-0.13,0.62) .. (0,0.72);
    \draw[c12flow] (0.2,2.55)--(0.2,2.05);
    \node[c12lab,right] at (0.22,2.35) {\ce{H2}};
    \node[c12lab] at (-0.22,1.0) {\ce{Cl2}};
    \node[c12lab] at (0,-0.6) {D．验证 \ce{Cl2} 的氧化性};
  \end{scope}
\end{tikzpicture}}{2024年广东卷第8题装置示意\label{fig:12-z24-8}}

\tcblower
\silu 逐项检查原理、试剂、装置。\ce{MnO2} 与盐酸反应制 \ce{Cl2}，要用浓盐酸并且加热，稀盐酸不反应；A 装置没有加热装置，不能达到目的。B：饱和食盐水吸收 HCl，同时 \ce{Cl-} 浓度大，抑制 \ce{Cl2} 溶解；无水 \ce{CaCl2} 是中性干燥剂，可干燥 \ce{Cl2}。C：\ce{Cl2} 密度比空气大，从长管进入，沉在瓶底，把空气从短管排出，即向上排空气法。D：先将氢气验纯、点燃，再把火焰伸入集氯气瓶。\ce{H2} 在 \ce{Cl2} 中安静地燃烧，发出苍白色火焰，\ce{H2} 被氧化，体现 \ce{Cl2} 的氧化性。\\
\daan A
\end{zhenti}

\section{物质的分离、提纯与检验}

分离是把混合物中的各成分分开，各自得到纯净物；提纯（除杂）只要求把杂质除去，得到被提纯的物质。选择方法的依据是各成分物理性质或化学性质的差异：颗粒大小不同用过滤，溶解度随温度变化不同用结晶，沸点不同用蒸馏，在两种溶剂中溶解度不同用萃取。

\subsection{过滤与蒸发}

过滤用于分离不溶性固体和液体。操作要点常概括为“一贴二低三靠”（图~\ref{fig:12-guolv}(a)）：滤纸紧贴漏斗内壁，中间不留气泡（一贴），否则过滤很慢；滤纸边缘低于漏斗边缘，液面低于滤纸边缘（二低），否则液体会从滤纸和漏斗之间流下，起不到过滤作用；倾倒液体的烧杯口紧靠玻璃棒，玻璃棒下端轻靠在三层滤纸一侧，漏斗下端管口紧靠烧杯内壁（三靠），分别防止液体洒出、防止戳破滤纸、防止滤液飞溅。

蒸发用于从溶液中得到溶解度受温度影响不大的固体溶质，例如从食盐水中得到 NaCl（图~\ref{fig:12-guolv}(b)）。加热时用玻璃棒不断搅拌，防止局部过热造成液滴飞溅；当蒸发皿中出现较多晶体时停止加热，利用余热把剩余的少量水蒸干。

\tu{\begin{tikzpicture}[scale=0.95]
  % (a) 过滤
  \fill[c12liq] (-0.59,0.02) rectangle (0.59,0.4);
  \draw[c12g] (-0.66,1.48)--(-0.6,1.42)--(-0.6,0)--(0.6,0)--(0.6,1.42)--(0.66,1.48);
  \fill[c12liq] (-0.96,2.6)--(-0.52,2.15)--(-0.44,2.15)--(0.0,2.6)--cycle;
  \draw[c12g] (-1.13,3.0)--(-0.52,2.12)--(-0.52,1.0) (0.17,3.0)--(-0.44,2.12)--(-0.44,1.05);
  \draw[c12g] (-0.52,1.0)--(-0.585,0.85)--(-0.44,1.05);
  \draw[c12g,dashed] (-1.02,2.88)--(-0.49,2.2)--(0.06,2.88);
  \draw[line width=1.1pt,black!65] (-0.62,2.33)--(0.45,3.9);
  \begin{scope}[shift={(-0.11,3.725)},rotate=-40]
    \draw[c12g] (-0.35,0)--(-0.35,1.0)--(0.35,1.0)--(0.35,0)--(0.43,-0.06);
    \fill[c12liq] (0.34,0.02)--(0.34,0.55)--(-0.34,0.0)--cycle;
  \end{scope}
  \draw[c12ptr] (-0.62,2.33)--(-1.6,2.2) node[left,c12lab]{玻璃棒下端靠在\\三层滤纸一侧};
  \draw[c12ptr] (-0.97,2.62)--(-1.6,3.0) node[left,c12lab]{液面低于滤纸边缘};
  \draw[c12ptr] (-0.58,0.9)--(-1.6,1.0) node[left,c12lab]{漏斗下端紧靠\\烧杯内壁};
  \draw[c12ptr] (0.16,3.55)--(0.75,3.95) node[right,c12lab]{烧杯口紧靠玻璃棒};
  \node[c12lab] at (-0.3,-0.4) {(a) 过滤};
  % (b) 蒸发
  \begin{scope}[xshift=5.6cm,yshift=1.4cm]
    \draw[line width=0.9pt,black!60] (0.6,-0.08)--(1.45,-0.08);
    \draw[line width=0.9pt,black!60] (1.45,-1.6)--(1.45,1.6);
    \fill[black!25] (-0.6,-1.68) rectangle (1.75,-1.6);
    \begin{scope}\clip (-0.72,0) arc[start angle=180,end angle=360,x radius=0.72,y radius=0.38]--cycle;
      \fill[c12liq] (-1,-0.5) rectangle (1,-0.12);\end{scope}
    \draw[c12g] (-0.78,0.03)--(-0.72,0) arc[start angle=180,end angle=360,x radius=0.72,y radius=0.38]--(0.72,0);
    \foreach \x/\y in {-0.3/-0.22,-0.1/-0.3,0.12/-0.26,0.3/-0.2} \fill[white] (\x,\y) rectangle ++(0.07,0.07);
    \draw[line width=0.9pt,black!60] (-0.62,-0.08)--(0.62,-0.08);
    \draw[line width=1.1pt,black!65] (-0.15,-0.25)--(0.55,1.2);
    \pic at (0,-0.24) {c12lamp};
    \draw[c12ptr] (0.45,1.0)--(-0.2,1.45) node[left,c12lab]{玻璃棒搅拌，防止\\局部过热液滴飞溅};
    \draw[c12ptr] (-0.55,-0.1)--(-1.15,0.35) node[left,c12lab]{蒸发皿};
    \draw[c12ptr] (1.45,0.6)--(2.0,0.85) node[right,c12lab]{铁架台\\（铁圈）};
    \node[c12lab] at (0.4,-1.95) {(b) 蒸发};
  \end{scope}
\end{tikzpicture}}{过滤与蒸发装置示意\label{fig:12-guolv}}

\subsection{蒸馏}

蒸馏利用各成分挥发性的差异进行分离。普通蒸馏适合沸点相差较大的互溶液体，也可除去液体中的难挥发杂质，例如由自来水制蒸馏水；它不能保证所有互溶液体都能一次分成纯物质。装置如图~\ref{fig:12-zhengliu}，蒸气的走向是：蒸馏烧瓶 $\to$ 支管 $\to$ 直形冷凝管 $\to$ 接液管（牛角管）$\to$ 锥形瓶。

\tu{\begin{tikzpicture}[scale=0.95]
  \begin{scope}\clip (0,0.72) circle (0.72);\fill[c12liq] (-1,-0.1) rectangle (1,0.6);\end{scope}
  \draw[c12g] (-0.17,2.65)--(-0.17,1.42) arc[start angle=103.66,end angle=436.34,radius=0.72]--(0.17,2.2);
  \draw[c12g] (0.17,2.65)--(0.17,2.3);
  \draw[c12t] (0.17,2.25)--($(0.17,2.25)+(-20:0.95)$);
  \draw[c12stop] (-0.21,2.6) rectangle (0.21,2.8);
  \draw[line width=0.4pt,fill=white] (-0.035,3.35) rectangle (0.035,2.32);
  \fill[black!70] (0,2.27) ellipse (0.05 and 0.08);
  \foreach \x/\y in {-0.3/0.12,0.05/0.08,0.32/0.15} \draw[line width=0.4pt,fill=white] (\x,\y) circle (0.06);
  \pic at (0,0) {c12net};
  \pic at (0,0) {c12lamp};
  % 冷凝管
  \begin{scope}[shift={($(0.17,2.25)+(-20:0.8)$)},rotate=-20]
    \fill[c12liq] (0.35,-0.2) rectangle (2.75,0.2);
    \draw[c12g] (0.35,-0.2) rectangle (2.75,0.2);
    \draw[c12t] (0,0)--(3.25,0);
    \draw[c12g] (2.5,-0.2)--(2.5,-0.48) (0.6,0.2)--(0.6,0.48);
    \draw[c12flow] (2.5,-0.82)--(2.5,-0.5);
    \draw[c12flow] (0.6,0.5)--(0.6,0.82);
    \coordinate (c12end) at (3.25,0);
    \coordinate (c12in) at (2.5,-0.85);
    \coordinate (c12out) at (0.6,0.85);
  \end{scope}
  \node[c12lab,right] at (c12in) {冷凝水下口进};
  \node[c12lab,above] at (c12out) {上口出};
  % 牛角管与锥形瓶
  \draw[c12t] (c12end)--++(-20:0.25)--++(-75:0.62);
  \pic at ($(c12end)+(-20:0.25)+(-75:0.62)+(0,-1.05)$) {c12cone};
  % 标注
  \draw[c12ptr] (0.05,2.27)--(-1.2,2.9) node[left,c12lab]{温度计水银球\\位于支管口处};
  \draw[c12ptr] (-0.3,0.12)--(-1.2,0.5) node[left,c12lab]{沸石（碎瓷片）\\防暴沸};
  \draw[c12ptr] (-0.5,-0.05)--(-1.2,-0.4) node[left,c12lab]{陶土网};
  \node[c12lab] at (1.5,-0.25) {蒸馏烧瓶内液体\\不超过容积的 2/3};
  \node[c12lab,below] at ($(c12end)+(0.6,-1.15)$) {接收器};
\end{tikzpicture}}{蒸馏装置示意\label{fig:12-zhengliu}}

读图时抓住三个位置。温度计测的是即将进入冷凝管的蒸气温度，所以水银球要放在蒸馏烧瓶支管口处；放得太低测到的是液体温度，太高则偏低，都不能准确判断馏分。烧瓶中加沸石或碎瓷片，它们表面多孔，能不断放出小气泡作为汽化中心，防止液体突然剧烈沸腾（暴沸）；如果加热后才发现忘加沸石，应停止加热，冷却后补加。冷凝水从冷凝管下口进、上口出，理由前面已经讲过。

\subsection{萃取与分液}

萃取是利用溶质在两种互不相溶的溶剂中溶解度的不同，用一种溶剂（萃取剂）把溶质从它与另一种溶剂组成的溶液中提取出来；分液是把两种互不相溶的液体分开。选择萃取剂要满足三点：与原溶剂互不相溶，不与溶质和原溶剂反应，溶质在萃取剂中的溶解度远大于在原溶剂中。从碘水中提取碘常用 \ce{CCl4}（密度比水大，萃取后在下层）或苯（密度比水小，在上层）；乙醇与水互溶，不能作萃取剂。

\tu{\begin{tikzpicture}[scale=0.95]
  % (a) 振荡与放气
  \begin{scope}[shift={(0,1.7)},rotate=-135]
    \pic[scale=1.1] {c12sepliq};
    \pic[scale=1.1] {c12sepbody};
    \draw[c12t] (0,-0.05)--(0,-0.75);
  \end{scope}
  \draw[c12flow] (-0.62,2.35)--(-0.95,2.68);
  \node[c12lab,above] at (-1.0,2.7) {放气};
  \node[c12lab] at (0.1,-0.35) {(a) 倒转振荡，\\打开活塞放气};
  % (b) 静置分液
  \begin{scope}[xshift=4.2cm]
    \draw[line width=0.9pt,black!60] (1.6,-0.1)--(1.6,4.3);
    \fill[black!25] (-0.8,-0.18) rectangle (1.9,-0.1);
    \draw[line width=0.9pt,black!60] (0.5,2.52)--(1.6,2.52);
    \draw[line width=0.9pt,black!60] (-0.42,2.52)--(0.42,2.52);
    \begin{scope}[shift={(0,2.05)},scale=1.3]
      \clip (0.04,0.07) .. controls (0.15,0.22) and (0.55,0.55) .. (0.5,0.95) .. controls (0.46,1.3) and (0.15,1.4) .. (0.1,1.45) -- (-0.1,1.45) .. controls (-0.15,1.4) and (-0.46,1.3) .. (-0.5,0.95) .. controls (-0.55,0.55) and (-0.15,0.22) .. (-0.04,0.07) -- cycle;
      \fill[exorange!30] (-0.6,0) rectangle (0.6,0.6);
      \fill[mainblue!12] (-0.6,0.6) rectangle (0.6,1.05);
      \draw[line width=0.3pt] (-0.6,0.6)--(0.6,0.6);
    \end{scope}
    \begin{scope}[shift={(0,2.05)},scale=1.3]
      \draw[c12g] (0.04,0.07) .. controls (0.15,0.22) and (0.55,0.55) .. (0.5,0.95) .. controls (0.46,1.3) and (0.15,1.4) .. (0.1,1.45) -- (0.1,1.6);
      \draw[c12g] (-0.04,0.07) .. controls (-0.15,0.22) and (-0.55,0.55) .. (-0.5,0.95) .. controls (-0.46,1.3) and (-0.15,1.4) .. (-0.1,1.45) -- (-0.1,1.6);
      \draw[c12g,fill=white] (-0.2,-0.045) rectangle (0.2,0.045);
      \draw[c12g,fill=white] (0.22,0) circle (0.055);
    \end{scope}
    \draw[c12stop] (-0.17,4.35) rectangle (0.17,4.57);
    \draw[c12flow] (0.3,4.2)--(0.3,4.55);
    \draw[c12t] (-0.4,2.0)--(-0.4,0.95);
    \draw[c12t] (0,1.98)--(0,2.0) ;
    \draw[c12g] (0.04,1.98)--(-0.36,1.98);
    \fill[exorange!30] (-0.59,0.02) rectangle (0.59,0.32);
    \draw[c12g] (-0.66,1.38)--(-0.6,1.32)--(-0.6,0)--(0.6,0)--(0.6,1.32)--(0.66,1.38);
    \draw[c12ptr] (0.45,3.0)--(-1.0,3.55) node[left,c12lab]{上层液体：\\从上口倒出};
    \draw[c12ptr] (-0.3,2.55)--(-1.0,2.4) node[left,c12lab]{下层液体：\\从下口放出};
    \draw[c12ptr] (-0.42,1.1)--(-1.0,1.0) node[left,c12lab]{漏斗下端紧靠\\烧杯内壁};
    \node[c12lab,right] at (0.35,4.75) {打开上口玻璃塞\\（或使凹槽对准小孔）};
    \node[c12lab] at (0.3,-0.45) {(b) 静置分层后分液};
  \end{scope}
\end{tikzpicture}}{萃取与分液操作示意\label{fig:12-fenye}}

操作过程见图~\ref{fig:12-fenye}。把溶液和萃取剂加入分液漏斗，塞好上口玻璃塞，倒转过来用力振荡，使溶质充分转移到萃取剂中；振荡过程中萃取剂挥发会使漏斗内压强增大，所以要不时打开活塞放气（图(a)）。然后把分液漏斗放在铁圈上静置，待液体分层。分液前先打开上口玻璃塞（或使塞上的凹槽对准漏斗口的小孔），使漏斗内外气压相通，否则液体流不下来。再打开活塞，下层液体从下口放出，待两层液体的界面恰好到达活塞时关闭活塞，上层液体从上口倒出（图(b)）。上层液体如果也从下口放出，就会被残留在活塞和下端管口中的下层液体污染。

\begin{liti}[基础]
用 \ce{CCl4} 从碘水中萃取碘。(1) 振荡、静置后，可观察到什么现象？(2) 能否用乙醇代替 \ce{CCl4}？(3) 分液后要从碘的 \ce{CCl4} 溶液中回收 \ce{CCl4}，应采用什么方法？
\tcblower
\jie (1) 液体分为两层，下层（\ce{CCl4} 层）呈紫红色，上层（水层）颜色变浅（接近无色）。\ce{CCl4} 密度比水大，所以在下层；\ce{I2} 在 \ce{CCl4} 中溶解度大，大部分转移到 \ce{CCl4} 层。\\
(2) 不能。乙醇与水互溶，加入后不分层，无法分液。\\
(3) 蒸馏。\ce{CCl4}（沸点约 77 ℃）与 \ce{I2} 沸点相差较大，蒸出的是 \ce{CCl4}。
\end{liti}

\subsection{结晶与重结晶}

从溶液中得到固体有两种结晶方式，选哪一种取决于溶质的溶解度随温度变化的大小。NaCl 的溶解度受温度影响很小，降温析出很少，通常用蒸发溶剂的办法，即蒸发结晶。\ce{KNO3}、苯甲酸的溶解度随温度升高明显增大，热的饱和溶液一冷却就析出大量晶体，适合冷却结晶，操作顺序是蒸发浓缩、冷却结晶、过滤、洗涤、干燥。带结晶水的晶体（如 \ce{CuSO4*5H2O}）和受热易分解的物质也要用冷却结晶，因为蒸干会使它们失去结晶水或分解。

重结晶用来提纯固体，例如提纯含少量 NaCl 和泥沙的苯甲酸粗品：加水加热使苯甲酸溶解，趁热过滤除去泥沙（趁热是为了防止苯甲酸在漏斗中析出而损失），滤液冷却结晶析出苯甲酸，过滤、洗涤得到较纯的晶体，而含量少的 NaCl 留在母液中。

除杂时要遵循四条原则：不引入新杂质，不减少被提纯的物质，杂质转化后容易分离，被提纯物质容易复原。常见的除杂方法如下（括号内为杂质）：\ce{CO2}（HCl）用饱和 \ce{NaHCO3} 溶液；\ce{Cl2}（HCl）用饱和食盐水；\ce{SO2}（HCl）用饱和 \ce{NaHSO3} 溶液；\ce{FeCl2} 溶液（\ce{FeCl3}）加过量铁粉后过滤；\ce{FeCl3} 溶液（\ce{FeCl2}）通 \ce{Cl2} 或加 \ce{H2O2}；\ce{NaHCO3} 溶液（\ce{Na2CO3}）通入过量 \ce{CO2}；乙酸乙酯（乙酸）加饱和 \ce{Na2CO3} 溶液后分液。除 \ce{CO2} 中的 HCl 不能用 NaOH 溶液或 \ce{Na2CO3} 溶液，因为它们也吸收 \ce{CO2}，违反“不减少被提纯物质”的原则。

\subsection{常见离子与气体的检验}

检验离子要写清“取样—加试剂—看现象—得结论”，并考虑其他离子的干扰。常见离子的检验方法归纳如下。

\begin{center}
\small
\begin{tabularx}{\linewidth}{p{1.4cm}X}
\toprule
离子 & 方法与现象 \\
\midrule
\ce{Fe^3+} & 加 KSCN 溶液，溶液变红；或加 NaOH 溶液生成红褐色沉淀 \\
\ce{Fe^2+} & 加 \ce{K3[Fe(CN)6]} 溶液生成蓝色沉淀（\ce{Fe^3+} 存在时也适用）；或先加 KSCN 溶液无现象，再加氯水变红（只适用于无 \ce{Fe^3+} 时） \\
\ce{NH4+} & 加浓 NaOH 溶液并加热，产生能使湿润的红色石蕊试纸变蓝的气体 \\
\ce{SO4^2-} & 先加足量稀盐酸，若有气体待其逸尽，若有沉淀先过滤，再向所得澄清液加 \ce{BaCl2}，生成白色沉淀（排除碳酸盐、亚硫酸盐和 \ce{Ag+} 等干扰） \\
\ce{Cl-} & 加稀硝酸酸化的 \ce{AgNO3} 溶液，生成白色沉淀 \\
\ce{I-} & 加少量新制氯水后再加淀粉溶液变蓝；或用 \ce{CCl4} 萃取，观察下层呈紫红色；氯水过量可能继续氧化碘 \\
\ce{Na+}、\ce{K+} & 焰色试验：Na 呈黄色；K 呈紫色（透过蓝色钴玻璃观察） \\
\bottomrule
\end{tabularx}
\end{center}

常见气体的检验：\ce{SO2} 使品红溶液褪色，加热后又恢复红色；\ce{CO2} 使澄清石灰水变浑浊（若同时有 \ce{SO2}，要先除去 \ce{SO2}，因为它也能使石灰水变浑浊）；\ce{NH3} 使湿润的红色石蕊试纸变蓝，遇蘸有浓盐酸的玻璃棒产生白烟；\ce{Cl2} 使湿润的淀粉-KI 试纸变蓝。

\section{气体的制备与性质实验}

\subsection{一套完整的气体实验装置}

制备一种气体并研究它的性质，一般要经过五个环节：在发生装置中生成气体，除去气体中的杂质，干燥，收集或进行性质实验，最后处理多余的有毒尾气。图~\ref{fig:12-cl2} 是实验室用 \ce{MnO2} 和浓盐酸制取并收集干燥、纯净 \ce{Cl2} 的装置，箭头表示气体流向。

\tu{\begin{tikzpicture}[scale=0.9]
  % 发生装置
  \pic at (0,0) {c12rbfliq};
  \pic at (0,0) {c12rbf};
  \pic at (0,0) {c12net};
  \pic at (0,0) {c12lamp};
  \draw[c12t] (-0.06,1.0)--(-0.06,2.55);
  \pic at (-0.06,2.6) {c12sepliq};
  \pic at (-0.06,2.6) {c12sepbody};
  \draw[c12ptr] (-0.35,3.25)--(-0.85,3.5) node[left,c12lab]{浓盐酸};
  \draw[c12ptr] (-0.4,0.3)--(-0.95,0.6) node[left,c12lab]{\ce{MnO2}};
  \node[c12lab] at (0,-1.75) {发生装置\\（固$+$液，加热）};
  % 导管
  \draw[c12t] (0.09,1.6)--(0.09,2.4)--(2.17,2.4);
  \draw[c12t] (2.43,2.4)--(3.97,2.4);
  \draw[c12t] (4.23,2.4)--(5.77,2.4);
  \draw[c12t] (6.03,2.4)--(7.57,2.4);
  \draw[c12t] (7.83,2.4)--(9.2,2.4)--(9.2,0.3);
  \draw[c12flow] (0.6,2.58)--(1.9,2.58);
  \node[c12lab,above] at (1.25,2.62) {\ce{Cl2}\\（含 HCl、\ce{H2O}）};
  \draw[c12flow] (2.65,2.58)--(3.75,2.58);
  \node[c12lab,above] at (3.2,2.62) {\ce{Cl2}\\（含 \ce{H2O}）};
  \draw[c12flow] (4.45,2.58)--(5.55,2.58);
  \node[c12lab,above] at (5.0,2.62) {干燥、纯净\\的 \ce{Cl2}};
  \draw[c12flow] (6.25,2.58)--(7.35,2.58);
  \node[c12lab,above] at (6.8,2.62) {多余 \ce{Cl2}};
  % 各装置
  \pic at (2.3,0) {c12wash};
  \node[c12lab] at (2.3,-0.5) {饱和食盐水\\[1pt]除杂：吸收 HCl};
  \pic at (4.1,0) {c12wash};
  \node[c12lab] at (4.1,-0.5) {浓硫酸\\[1pt]干燥};
  \fill[c12gas] (5.46,0.04) rectangle (6.34,0.7);
  \pic at (5.9,0) {c12bottle};
  \node[c12lab] at (5.9,-0.5) {向上排空气法\\[1pt]收集（长进短出）};
  \pic at (7.7,0) {c12safe};
  \node[c12lab] at (7.7,-0.5) {空安全瓶\\[1pt]缓冲倒吸};
  \pic at (9.2,0) {c12beaker};
  \node[c12lab] at (9.2,-0.5) {NaOH 溶液\\[1pt]尾气处理};
\end{tikzpicture}}{实验室制取干燥、纯净的 \ce{Cl2}（装置连接示意）\label{fig:12-cl2}}

按气流方向读这张图。烧瓶中发生的反应是
\[\ce{MnO2 + 4HCl}\text{（浓）}\ce{->[\Delta] MnCl2 + Cl2 ^ + 2H2O}\]
浓盐酸易挥发，加热时又有水蒸气，所以出来的 \ce{Cl2} 中混有 HCl 和 \ce{H2O}。第一个洗气瓶盛饱和食盐水，HCl 极易溶于水而被吸收，饱和食盐水中 \ce{Cl-} 浓度大，使 \ce{Cl2 + H2O <=> H+ + Cl- + HClO} 的平衡左移，\ce{Cl2} 损失很少。第二个洗气瓶盛浓硫酸，吸收水蒸气。干燥的 \ce{Cl2} 进入集气瓶，它的密度比空气大，从长管进入后沉在瓶底，空气从短管排出，这就是向上排空气法；当集气瓶上方充满黄绿色时，说明已收集满。最后多余的 \ce{Cl2} 通入 NaOH 溶液：\ce{Cl2 + 2NaOH = NaCl + NaClO + H2O}。

连接顺序的依据有三条。第一，先除杂后干燥：水溶液洗气会重新带出水蒸气。第二，洗气瓶长管进、短管出，使气体通过液体；接反会把液体从长管压出。收集瓶图示为正放的双导管瓶：密度比空气大的气体长进短出，密度小的短进长出；瓶若倒放，必须重新按气体实际进入的位置判断。第三，尾气吸收放在最后。吸收反应很快或停止加热后气体冷却，均可能使压强骤降；图中空安全瓶的两管均不触液，倒吸液先留在瓶内。安全瓶容积须能容纳可能倒吸的液体，它不能替代气密性检查和规范操作。

\subsection{常见气体的实验室制法}

选择发生装置看反应物的状态和反应条件。固体与固体加热，如用 \ce{NH4Cl} 和 \ce{Ca(OH)2} 制 \ce{NH3}，用试管，管口略向下倾斜（防止生成的水倒流使试管炸裂）；固体与液体不加热，如制 \ce{H2}、\ce{CO2}、\ce{SO2}，以及用 \ce{KMnO4} 与浓盐酸制 \ce{Cl2}；固体与液体加热，如 \ce{MnO2} 与浓盐酸制 \ce{Cl2}。选择干燥剂看是否与气体反应：浓硫酸是酸性、强氧化性干燥剂，不能干燥 \ce{NH3}、\ce{H2S}；碱石灰是碱性干燥剂，不能干燥 \ce{Cl2}、\ce{SO2}、\ce{CO2} 等酸性气体；无水 \ce{CaCl2} 是中性干燥剂，但能与 \ce{NH3} 结合，也不能干燥 \ce{NH3}。

\begin{center}
\small
\begin{tabularx}{\linewidth}{p{1.0cm}p{6.6cm}X}
\toprule
气体 & 实验室制法 & 净化、收集、尾气处理 \\
\midrule
\ce{Cl2} & \ce{MnO2 + 4HCl}（浓）\ce{->[\Delta] MnCl2 + Cl2 ^ + 2H2O}；或 \ce{2KMnO4 + 16HCl}（浓）\ce{= 2KCl + 2MnCl2 + 5Cl2 ^ + 8H2O} & 饱和食盐水除 HCl，浓硫酸干燥；向上排空气法；NaOH 溶液吸收 \\
\ce{NH3} & \ce{2NH4Cl + Ca(OH)2 ->[\Delta] CaCl2 + 2NH3 ^ + 2H2O}；或把浓氨水滴到 NaOH 固体或 CaO 上 & 碱石灰干燥；向下排空气法；水或稀硫酸吸收（防倒吸） \\
\ce{SO2} & \ce{Na2SO3 + H2SO4}（浓）\ce{= Na2SO4 + SO2 ^ + H2O}（常用 70\% 左右的硫酸） & 浓硫酸干燥；向上排空气法；NaOH 溶液吸收 \\
\ce{CO2} & \ce{CaCO3 + 2HCl = CaCl2 + CO2 ^ + H2O} & 饱和 \ce{NaHCO3} 溶液除 HCl，浓硫酸干燥；向上排空气法 \\
\bottomrule
\end{tabularx}
\end{center}

\begin{liti}[基础]
实验室用大理石与稀盐酸反应，制取并收集一瓶干燥、纯净的 \ce{CO2}。可选用的装置有：①盛浓硫酸的洗气瓶；②盛饱和 \ce{NaHCO3} 溶液的洗气瓶；③盛 NaOH 溶液的洗气瓶；④用向上排空气法收集的集气瓶。(1) 生成的 \ce{CO2} 中含有哪些杂质？(2) 选择装置并按气流方向排序。(3) 为什么不选③？
\tcblower
\jie (1) 盐酸挥发出的 HCl 和水蒸气。\\
(2) 发生装置 $\to$ ② $\to$ ① $\to$ ④。先通过饱和 \ce{NaHCO3} 溶液除去 HCl（\ce{NaHCO3 + HCl = NaCl + CO2 ^ + H2O}），再通过浓硫酸干燥；顺序不能颠倒，否则气体从 \ce{NaHCO3} 溶液中又带出水蒸气。\ce{CO2} 密度比空气大，用向上排空气法收集。\\
(3) NaOH 溶液同时吸收 \ce{CO2}（\ce{CO2 + 2NaOH = Na2CO3 + H2O}），会使被提纯的气体减少。
\end{liti}

\begin{zhenti}{改编自2026年广东卷第9题}
用图~\ref{fig:12-z26-9} 所示装置进行实验：\ce{Ca(ClO)2} 中滴加盐酸产生 \ce{Cl2}，气体依次通过 KI 溶液、饱和食盐水（瓶上插有长颈漏斗）、浓硫酸，进入装有铁丝的 U 形管（加热），最后用 NaOH 溶液吸收尾气。反应一段时间后，关闭活塞和止水夹，熄灭酒精灯，冷却后提起铁丝，注入稀硫酸。下列说法错误的是\xz
\tu{\begin{tikzpicture}[scale=0.82]
  \fill[black!70] (-0.3,0.12) circle (0.05) (-0.1,0.08) circle (0.05) (0.12,0.1) circle (0.05) (0.3,0.15) circle (0.05);
  \pic at (0,0) {c12rbf};
  \draw[c12t] (-0.06,0.9)--(-0.06,2.55);
  \pic at (-0.06,2.6) {c12sepliq};
  \pic at (-0.06,2.6) {c12sepbody};
  \draw[c12ptr] (-0.35,3.25)--(-0.8,3.5) node[left,c12lab]{盐酸};
  \draw[c12ptr] (-0.3,0.15)--(-0.8,0.5) node[left,c12lab]{\ce{Ca(ClO)2}};
  \node[c12lab,left] at (-0.12,2.85) {活塞};
  \draw[c12t] (0.09,1.6)--(0.09,2.4)--(1.77,2.4);
  \draw[c12t] (2.03,2.4)--(3.47,2.4);
  \draw[c12t] (3.73,2.4)--(5.17,2.4);
  \draw[c12flow] (0.55,2.58)--(1.5,2.58);
  \pic at (1.9,0) {c12wash};
  \node[c12lab] at (1.9,-0.25) {KI 溶液};
  \pic at (3.6,0) {c12wash};
  \draw[c12t] (3.6,0.25)--(3.6,3.0);
  \draw[c12g] (3.57,3.0)--(3.33,3.38) (3.63,3.0)--(3.87,3.38);
  \node[c12lab] at (3.6,-0.25) {饱和食盐水};
  \pic at (5.3,0) {c12wash};
  \node[c12lab] at (5.3,-0.25) {浓硫酸};
  % U 形管
  \draw[c12t] (5.43,2.4)--(6.7,2.4)--(6.7,1.4);
  \draw[c12g] (6.6,1.65)--(6.6,0.5) arc[start angle=180,end angle=360,radius=0.5]--(7.6,1.65);
  \draw[c12g] (6.95,1.65)--(6.95,0.5) arc[start angle=180,end angle=360,radius=0.15]--(7.25,1.65);
  \draw[c12stop] (6.55,1.6) rectangle (7.0,1.78);
  \draw[c12stop] (7.2,1.6) rectangle (7.65,1.78);
  \draw[c12wire] (6.85,2.15)--(6.85,1.3)--(6.72,1.2)--(6.9,1.05)--(6.72,0.9)--(6.9,0.75)--(6.72,0.6)--(6.85,0.5);
  \draw[c12wire] (6.85,2.15) circle (0.06);
  \draw[c12ptr] (6.9,1.0)--(6.2,0.6) node[left,c12lab]{铁丝};
  \pic at (7.1,0) {c12lamp};
  \draw[c12t] (7.43,1.45)--(7.43,2.4)--(8.77,2.4);
  \draw[c12t] (9.03,2.4)--(10.4,2.4)--(10.4,0.3);
  \draw[line width=1.3pt] (8.12,2.22)--(8.12,2.58) (8.24,2.22)--(8.24,2.58);
  \node[c12lab,above] at (8.18,2.6) {止水夹};
  \pic at (8.9,0) {c12safe};
  \node[c12lab] at (8.9,-0.25) {空安全瓶};
  \pic at (10.4,0) {c12beaker};
  \node[c12lab] at (10.4,-0.25) {NaOH 溶液};
\end{tikzpicture}}{氯气性质实验与尾气吸收装置\label{fig:12-z26-9}}
\xxb{KI 溶液变棕黄色，体现 \ce{I-} 的还原性}{浓硫酸的作用是吸收水分，干燥 \ce{Cl2}}{U 形管中溶液遇 KSCN 溶液变红，说明有 \ce{Fe^3+} 存在}{烧杯溶液中最终存在 $c(\ce{Na+})+c(\ce{H+})=c(\ce{Cl-})+c(\ce{ClO-})$}
\tcblower
\silu 装置中的反应是 \ce{Ca(ClO)2 + 4HCl = CaCl2 + 2Cl2 ^ + 2H2O}，不需加热。饱和食盐水瓶上的长颈漏斗起平衡气压的作用：若后面装置堵塞，瓶内压强增大，漏斗中液面上升，起到安全瓶的作用。\\
A：\ce{Cl2 + 2I- = 2Cl- + I2}，\ce{I-} 被氧化，体现还原性，正确。B 正确，干燥是为了防止水蒸气进入加热的 U 形管。C：Fe 在 \ce{Cl2} 中燃烧生成 \ce{FeCl3}，溶于稀硫酸后遇 KSCN 变红，正确。D：NaOH 吸收 \ce{Cl2} 后的溶液中还有 \ce{OH-}，电荷守恒应为 $c(\ce{Na+})+c(\ce{H+})=c(\ce{Cl-})+c(\ce{ClO-})+c(\ce{OH-})$，错误。\\
\daan D
\end{zhenti}

\begin{liti}[综合·易潮解物质的制备]
实验室用 \ce{Cl2} 与铁粉加热制备无水 \ce{FeCl3}（\ce{FeCl3} 易潮解，受热易升华）。装置依次为：A（\ce{MnO2} 与浓盐酸加热制 \ce{Cl2}）$\to$ B $\to$ C $\to$ D（硬质玻璃管中加热铁粉）$\to$ E（广口瓶，冷水浴收集产品）$\to$ F $\to$ G（带防倒吸措施的 NaOH 尾气吸收装置）。(1) B、C 中分别盛放什么试剂？(2) 为什么要先通一段时间 \ce{Cl2} 再点燃 D 处酒精灯？(3) 产品为什么在 E 中而不是 D 中得到？(4) F 中应盛放什么？作用是什么？
\tcblower
\silu 前半部分 A $\to$ B $\to$ C 与图~\ref{fig:12-cl2} 相同；这里的“收集”换成了制备产品的 D、E，后面多出的 F 是本题的关键。\\
\jie (1) B：饱和食盐水（除 HCl）；C：浓硫酸（干燥）。\\
(2) 排尽装置中的空气，防止铁被 \ce{O2} 氧化，使产品混有杂质。\\
(3) 生成的 \ce{FeCl3} 受热升华，随气流进入 E 中，遇冷凝华为固体。\\
(4) 浓硫酸（或装有无水 \ce{CaCl2} 的干燥管），防止 G 中的水蒸气进入 E 使 \ce{FeCl3} 潮解。\\
\dianbo 制备易水解、易潮解的物质，产品前后都要防水：前面干燥气体，后面加干燥装置，隔绝尾气处理装置中的水蒸气。
\end{liti}

\section{定量实验}\label{sec:12-dingliang}

定量实验要测出一个量（浓度、纯度、组成），关键是两件事：测量要准确，计算要找对守恒关系。中学常用滴定法、气体体积法和重量法。分析误差时，先写出计算待测量的式子，再看每个错误操作使式中哪个测量值偏大或偏小。

\subsection{滴定法}

中和滴定用已知浓度的标准溶液滴定未知浓度的待测溶液，根据恰好反应时消耗标准溶液的体积求待测溶液浓度。对一元酸与一元碱，
\[c(\text{待测})=\frac{c(\text{标准})\cdot V(\text{标准})}{V(\text{待测})}\]
本式适用于 $1:1$ 的反应；其他计量比应先写方程式确定。分析常见操作误差时，先把计算式中采用的 $c(\text{标准})$ 和 $V(\text{待测})$ 看作记录值，再判断标准液的读数差怎样改变：读数差偏大，计算结果偏高；偏小则偏低。不能脱离计算式，把任何实验的误差都机械归结为同一个体积。

\tu{\begin{tikzpicture}[scale=0.95]
  % (a) 滴定装置
  \fill[black!25] (-1.6,-0.1) rectangle (1.0,0);
  \draw[line width=1pt,black!60] (-1.1,0)--(-1.1,5.3);
  \draw[line width=1pt,black!60] (-1.1,3.7)--(-0.12,3.7);
  \draw[line width=0.6pt] (-0.15,3.6) rectangle (-0.09,3.8) (0.09,3.6) rectangle (0.15,3.8);
  \fill[c12liq] (-0.09,1.95) rectangle (0.09,4.6);
  \draw[c12g] (-0.09,5.2)--(-0.09,1.95) (0.09,5.2)--(0.09,1.95);
  \foreach \i in {0,...,7} \draw[line width=0.3pt] (-0.09,4.95-0.4*\i)--(0.0,4.95-0.4*\i);
  \draw[c12g,fill=white] (0,1.85) circle (0.1);
  \draw[line width=1pt] (-0.1,1.85)--(-0.32,1.85);
  \draw[c12g] (-0.04,1.76)--(-0.015,1.15) (0.04,1.76)--(0.015,1.15);
  \fill[white,draw=black!40,line width=0.3pt] (-0.85,0) rectangle (0.85,0.04);
  \pic[c12coneliq/.style={fill=red!12}] at (0,0.05) {c12cone};
  \draw[c12ptr] (-0.32,1.85)--(-1.4,2.2) node[left,c12lab]{左手控制活塞};
  \draw[c12ptr] (0.5,0.3)--(1.3,0.6) node[right,c12lab]{右手摇动锥形瓶\\眼睛注视瓶内\\溶液颜色变化};
  \draw[c12ptr] (0.75,0.02)--(1.3,-0.25) node[right,c12lab]{下垫白纸};
  \draw[c12ptr] (0.09,4.3)--(0.7,4.6) node[right,c12lab]{滴定管};
  \node[c12lab] at (-0.2,-0.45) {(a) 滴定操作};
  % (b) 读数
  \begin{scope}[xshift=5.6cm,yshift=1.0cm]
    \fill[c12liq] (0,0) .. controls (0.15,-0.12) and (0.35,-0.12) .. (0.5,0)--(0.5,-1.2)--(0,-1.2)--cycle;
    \draw[c12g] (0,1.4)--(0,-1.2) (0.5,1.4)--(0.5,-1.2);
    \draw[line width=0.4pt] (0,0) .. controls (0.15,-0.12) and (0.35,-0.12) .. (0.5,0);
    \foreach \i/\t in {0/20,1/21,2/22} {\draw[line width=0.4pt] (0,1.0-0.9*\i)--(0.2,1.0-0.9*\i); \node[font=\tiny,right] at (0.5,1.0-0.9*\i) {\t};}
    \foreach \i in {0,...,8} \draw[line width=0.25pt] (0,0.91-0.09*\i)--(0.1,0.91-0.09*\i);
    \foreach \i in {0,...,8} \draw[line width=0.25pt] (0,0.01-0.09*\i)--(0.1,0.01-0.09*\i);
    \coordinate (c12m) at (0.25,-0.09);
    \draw[line width=0.45pt,dashed] (-1.9,0.75)--(c12m);
    \draw[line width=0.45pt] (-1.9,-0.09)--(c12m);
    \draw[line width=0.45pt,dashed] (-1.9,-0.93)--(c12m);
    \node[c12lab,left] at (-1.9,0.75) {俯视：读数偏小};
    \node[c12lab,left] at (-1.9,-0.09) {平视：正确};
    \node[c12lab,left] at (-1.9,-0.93) {仰视：读数偏大};
    \node[c12lab] at (-0.8,-1.75) {(b) 滴定管读数（刻度从上往下增大）};
  \end{scope}
\end{tikzpicture}}{滴定操作与滴定管读数示意\label{fig:12-diding}}

图~\ref{fig:12-diding}(a) 是滴定操作的示意。完整的步骤是：检查滴定管是否漏水；用蒸馏水洗净后，再用待装溶液润洗 2～3 次；装液，赶出尖嘴部分的气泡，调节液面到“0”或“0”以下某一刻度，记下初读数；用滴定管（或移液管）准确量取待测液放入锥形瓶，加 2～3 滴指示剂；左手控制活塞，右手摇动锥形瓶，眼睛注视锥形瓶内溶液颜色的变化，锥形瓶下垫白纸便于观察；当滴入最后半滴标准溶液，溶液颜色发生突变且半分钟内不恢复时，即为滴定终点，记下终读数；重复 2～3 次，取平均值。

图~\ref{fig:12-diding}(b) 说明了读数误差的来源。视线与凹液面最低处相平才是正确读数。滴定管的刻度在管的前壁上，俯视时视线经过凹液面最低处，与前壁的交点在液面上方，而滴定管的刻度从上往下增大，所以读数偏小；仰视则读数偏大。量筒的刻度从下往上增大，俯视读数反而偏大，两者不要混淆。

\begin{center}
\small
\begin{tabularx}{\linewidth}{Xp{3.4cm}p{1.6cm}}
\toprule
操作 & 对 $V(\text{标准})$ 的影响 & 测定结果 \\
\midrule
酸式滴定管水洗后未用标准溶液润洗（标准溶液被稀释） & 偏大 & 偏高 \\
锥形瓶用待测液润洗（待测物增多） & 偏大 & 偏高 \\
锥形瓶洗净后残留少量蒸馏水 & 不变 & 无影响 \\
滴定前尖嘴有气泡，滴定后气泡消失 & 读数差偏大 & 偏高 \\
滴定前平视读数，滴定后俯视读数 & 终读数偏小 & 偏低 \\
量取待测液的滴定管未用待测液润洗（待测物减少） & 偏小 & 偏低 \\
摇动时锥形瓶中溶液溅出 & 偏小 & 偏低 \\
\bottomrule
\end{tabularx}
\end{center}

氧化还原滴定的操作与中和滴定相同，只是反应换成了氧化还原反应。常见的有高锰酸钾法（\ce{MnO4^- + 5Fe^2+ + 8H+ = Mn^2+ + 5Fe^3+ + 4H2O}，\ce{KMnO4} 自身紫红色，不另加指示剂）和碘量法（\ce{I2 + 2S2O3^2- = 2I- + S4O6^2-}，用淀粉作指示剂）。此外还有沉淀滴定和配位滴定，具体试剂和终点判据应由题目给出。计算时用关系式法：由方程式或电子守恒找出待测物与滴定剂之间物质的量的比例关系，再注意取样比例（如从 100.00 mL 中取 25.00 mL）。

\begin{liti}[基础]
用 0.1000\molL{} 盐酸滴定 20.00 mL 未知浓度的 NaOH 溶液，以酚酞为指示剂，三次实验平均消耗盐酸 22.00 mL。(1) 求 NaOH 溶液的浓度。(2) 如何判断滴定终点？(3) 若盛盐酸的滴定管滴定前尖嘴有气泡、滴定后消失，测得的 NaOH 浓度偏高还是偏低？
\tcblower
\jie (1) $c(\ce{NaOH})=\dfrac{0.1000\times22.00}{20.00}=0.1100$\molL。\\
(2) 滴入最后半滴盐酸时，溶液由浅红色变为无色，且半分钟内不恢复红色。\\
(3) 滴定前气泡占据了一部分体积，滴定后气泡被液体填充，这部分体积也被读进了“消耗的盐酸”中，$V(\text{标准})$ 偏大，结果偏高。
\end{liti}

\begin{liti}[典型·氧化还原滴定]
称取 2.000 g 某 \ce{FeSO4*7H2O}（$M=278$ g·mol$^{-1}$）样品（杂质不参与滴定反应），溶于稀硫酸配成 100.00 mL 溶液。取 25.00 mL，用 0.02000\molL{} \ce{KMnO4} 标准溶液滴定，消耗 16.00 mL。(1) 判断滴定终点的现象。(2) 求样品中 \ce{FeSO4*7H2O} 的质量分数。(3) \ce{KMnO4} 溶液应装在哪种滴定管中？
\tcblower
\jie (1) 滴入最后半滴 \ce{KMnO4} 溶液时，溶液出现过量高锰酸根的浅紫红色，且半分钟内不褪色；不能要求含 \ce{Fe^3+} 的待测液原本完全无色。\\
(2) 由 \ce{MnO4^-} $\sim$ 5\ce{Fe^2+}：25.00 mL 中 $n(\ce{Fe^2+})=5\times0.02000\times0.01600=1.600\times10^{-3}$ mol，100.00 mL 中共 $6.400\times10^{-3}$ mol，质量 $6.400\times10^{-3}\times278=1.779$ g，质量分数 $\dfrac{1.779}{2.000}\times100\%\approx88.96\%$。\\
(3) 酸式滴定管（或聚四氟乙烯活塞滴定管）。\ce{KMnO4} 有强氧化性，会腐蚀碱式滴定管的橡胶管。
\end{liti}

\subsection{气体体积法}

有气体生成的反应，可以通过测量气体体积来计算物质的量，常用图~\ref{fig:12-liangqi} 的量气装置。量气管（可用碱式滴定管改装）和水准管下端用橡胶管连通，构成连通器。反应产生的气体进入量气管上部，把水压下去，前后两次读数之差就是气体体积。

\tu{\begin{tikzpicture}[scale=0.95]
  \pic at (0,0) {c12cone};
  \foreach \x in {-0.3,-0.05,0.2} \fill[black!60] (\x,0.06) rectangle ++(0.14,0.04);
  \draw[c12stop] (-0.19,1.25) rectangle (0.19,1.42);
  % 恒压滴液漏斗
  \fill[c12liq] (-0.17,2.05) rectangle (0.07,2.6);
  \fill[c12liq] (-0.17,2.05)--(-0.08,1.9)--(-0.02,1.9)--(0.07,2.05)--cycle;
  \draw[c12g] (-0.17,3.15)--(-0.17,2.05)--(-0.08,1.9)--(-0.08,1.85) (0.07,3.15)--(0.07,2.05)--(-0.02,1.9)--(-0.02,1.85);
  \draw[c12stop] (-0.2,3.12) rectangle (0.1,3.27);
  \draw[c12g,fill=white] (-0.2,1.78) rectangle (0.1,1.85);
  \draw[c12t] (-0.05,1.77)--(-0.05,0.95);
  \draw[c12t,mainblue!70!black] (0.07,3.0)--(0.35,3.0)--(0.35,1.55)--(-0.13,1.55)--(-0.13,1.05);
  \draw[c12ptr] (-0.17,2.4)--(-0.75,2.7) node[left,c12lab]{稀盐酸\\（恒压滴液漏斗）};
  \draw[c12ptr] (-0.25,0.1)--(-0.8,0.4) node[left,c12lab]{镁带};
  % 导管到量气管
  \draw[c12t] (0.11,1.3)--(0.11,1.42);
  \draw[c12t] (0.11,1.42)--(0.11,1.5) -- (0.55,1.5)--(0.55,4.0)--(2.4,4.0)--(2.4,3.75);
  \draw[c12flow] (1.0,4.17)--(1.9,4.17);
  % 量气管
  \fill[c12liq] (2.25,0.6) rectangle (2.55,2.2);
  \draw[c12g] (2.25,3.75) rectangle (2.55,0.6);
  \draw[c12g,fill=white] (2.3,3.75) rectangle (2.5,3.85);
  \foreach \i in {0,...,7} \draw[line width=0.3pt] (2.25,3.45-0.38*\i)--(2.37,3.45-0.38*\i);
  \node[font=\tiny,left] at (2.25,3.45) {0};
  % 水准管
  \fill[c12liq] (3.55,1.0) rectangle (3.85,2.2);
  \draw[c12g] (3.55,3.2)--(3.55,1.0) (3.85,3.2)--(3.85,1.0);
  % 橡胶管
  \draw[line width=2.2pt,black!55] (2.4,0.6)--(2.4,0.25) .. controls (2.4,0.0) and (3.7,0.0) .. (3.7,0.25)--(3.7,1.0);
  \draw[c12g,dashed] (2.0,2.2)--(4.2,2.2);
  \draw[c12flow,<->] (4.1,2.6)--(4.1,3.2);
  \node[c12lab,right] at (4.15,2.9) {上下移动};
  \node[c12lab] at (2.4,-0.3) {量气管};
  \node[c12lab] at (3.7,-0.3) {水准管};
  \node[c12lab,right] at (4.2,2.2) {两液面相平};
\end{tikzpicture}}{测量气体体积的装置示意\label{fig:12-liangqi}}

读数前要做好三件事。第一，冷却到室温：反应放热，气体受热膨胀，趁热读数体积偏大。第二，上下移动水准管，使两管液面相平：此时量气管内气体压强等于大气压，测得的才是该温度和大气压下的体积。如果水准管液面高于量气管液面，管内气体受到额外的水柱压强而被压缩，读出的体积偏小。第三，读数时视线与凹液面最低处相平。另外，图中用恒压滴液漏斗加酸：如果改用普通分液漏斗，加入的盐酸会占据瓶中一部分空间，把等体积的空气挤进量气管，测得的气体体积偏大；恒压滴液漏斗的侧管使漏斗上方与锥形瓶连通，液体流下时总体积不变，避免了这一误差。

\begin{liti}[典型·气体体积法测定纯度]
用图~\ref{fig:12-liangqi} 所示原理、选择量程足够的量气管，测定某镁带样品中 Mg 的质量分数（杂质不与酸反应）。取样品 0.120 g，实验结束后测得生成气体的体积为 100.8 mL（已扣除水蒸气影响并换算为标准状况）。(1) 求 Mg 的质量分数。(2) 若读数时水准管液面高于量气管液面，测得结果偏大还是偏小？(3) 若反应结束后未冷却就立即读数，结果如何？
\tcblower
\jie (1) \ce{Mg + 2HCl = MgCl2 + H2 ^}，$n(\ce{H2})=\dfrac{0.1008}{22.4}=4.5\times10^{-3}$ mol $=n(\ce{Mg})$，$m(\ce{Mg})=0.108$ g，质量分数为 90.0\%。\\
(2) 水准管液面高，量气管内气体压强大于大气压，气体被压缩，读出的体积偏小，结果偏小。\\
(3) 气体温度高于室温，体积偏大，结果偏大。
\end{liti}

重量法把待测成分转化为组成确定的沉淀，或用吸收剂吸收生成的气体，通过质量求含量。沉淀要完全生成，并洗净、干燥或灼烧至恒重：重复“加热、在干燥器中冷却、称量”，相邻两次质量差达到实验规定的允许范围，不能一律规定为 0.1 g。称量前必须冷却，热容器周围的气流会干扰读数。若测吸收装置的增重，还要防空气中的 \ce{CO2} 和水蒸气进入，并用经过适当净化的载气把残留的待测气体全部赶入吸收装置。

\section{探究性实验与方案设计}

\subsection{控制变量与对照}

一个现象往往受几个因素影响，例如反应速率与浓度、温度、催化剂都有关。要弄清某一个因素的作用，就让其他因素保持不变，只改变这一个因素，观察结果的变化，这种方法叫控制变量法。把改变了该因素的实验和没有改变的实验放在一起比较，后者就是对照实验。设计这类表格时，常用的技巧是保持溶液总体积相同：改变某种试剂的体积时用水补足，这样其他物质的浓度才不变。

\begin{liti}[基础·控制变量]
用 \ce{Na2S2O3} 溶液与稀硫酸的反应（\ce{Na2S2O3 + H2SO4 = Na2SO4 + S v + SO2 ^ + H2O}）探究外界条件对反应速率的影响，实验记录如下（溶液浓度均为 0.10\molL）。

\begin{center}
\small
\begin{tabular}{ccccc}
\toprule
实验 & 温度/℃ & $V(\ce{Na2S2O3})$/mL & $V(\ce{H2SO4})$/mL & $V(\ce{H2O})$/mL \\
\midrule
① & 25 & 10.0 & 10.0 & 0 \\
② & 25 & 5.0 & 10.0 & $x$ \\
③ & 35 & 10.0 & 10.0 & 0 \\
\bottomrule
\end{tabular}
\end{center}
(1) $x$ 为多少？实验①②探究的是什么因素？(2) 探究温度的影响应比较哪两组实验？(3) 用什么现象比较反应的快慢？
\tcblower
\jie (1) $x=5.0$。三组溶液总体积都应为 20.0 mL，这样②中 \ce{H2SO4} 的浓度与①相同，只有 \ce{Na2S2O3} 的浓度不同，所以①②探究 \ce{Na2S2O3} 浓度对反应速率的影响。若②不加水，\ce{H2SO4} 浓度也变了，就无法判断速率变化是哪个因素引起的。\\
(2) ①和③，两者只有温度不同。\\
(3) 生成的 S 不溶于水，溶液变浑浊。可在锥形瓶底下的白纸上画一个“十”字，记录从混合到看不清“十”字所需的时间，时间越短，反应越快。
\end{liti}

\subsection{提出假设与验证}

探究实验一般按这样的思路进行：提出问题，作出假设（往往有两种或多种可能），设计实验（控制变量，设置对照或空白实验），实验并记录，分析得出结论，反思评价。假设要有依据，每一种假设都要找到能把它和其他假设区分开的实验现象。

高考题常常从“结果与预期不符”出发。2024年广东卷第17题中，测得体系 pH 不再变化，甲同学认为反应已达平衡，乙同学认为证据不足，提出“硫酸钙已耗尽”和“硫酸钙有剩余但被碳酸钙包裹”两种假设，再设计实验逐一检验；2026年广东卷第17题中，实验测得的两段消耗酸的体积之比与理论预测不符，进而提出“沉淀中还含有 \ce{Cu2(OH)2SO4}”的假设并定量验证。

空白实验是在不加待测物、其他条件完全相同的情况下进行的实验，用来扣除试剂、溶剂本身带来的影响。2026年广东卷第17题中，先用同浓度硫酸滴定等体积的稀氨水，测出氨水自身消耗的酸 $V_1$，才能由第一个终点的总体积 $V_2$ 求出铜氨配合物反应消耗的酸 $V_2-V_1$。

\begin{liti}[综合·反应限度的探究]
某小组探究 \ce{FeCl3} 溶液与 KI 溶液的反应 \ce{2Fe^3+ + 2I- <=> 2Fe^2+ + I2} 是否存在限度。实验：取 5 mL 0.1\molL{} KI 溶液，滴加 1 mL 0.1\molL{} \ce{FeCl3} 溶液，充分反应。(1) 加入 \ce{CCl4} 振荡、静置，下层呈紫红色，说明什么？(2) 取分液后的上层溶液，应加入什么试剂、出现什么现象，才能说明该反应存在限度？(3) 为什么实验中要让 KI 过量，而不是让 \ce{FeCl3} 过量？(4) 有同学提出：检验反应后溶液中仍有 \ce{I-}，也能说明反应存在限度。这种说法是否正确？
\tcblower
\silu “存在限度”意味着两种反应物都不能完全转化。要证明这一点，应让一种反应物过量，再检验另一种（不足量的）反应物是否仍有剩余。\\
\jie (1) 有 \ce{I2} 生成，即反应确实发生了。\\
(2) 加入 KSCN 溶液，溶液变红，说明 KI 过量时仍有 \ce{Fe^3+} 剩余，\ce{Fe^3+} 没有反应完全，反应存在限度。\\
(3) $n(\ce{I-})=5\times10^{-4}$ mol，$n(\ce{Fe^3+})=1\times10^{-4}$ mol，按方程式 \ce{I-} 远远过量。若 \ce{Fe^3+} 过量，反应后有 \ce{Fe^3+} 剩余是必然的，不能说明问题。\\
(4) 不正确。KI 本来就过量，即使反应能进行到底，溶液中也有 \ce{I-}。检验的对象必须是不足量的反应物。
\end{liti}

\subsection{方案的规范表述}

答题时要把操作、现象和结论写完整。描述现象要写看到的（颜色、沉淀、气体），不能写结论，例如写“生成红褐色沉淀”，而不是“生成 \ce{Fe(OH)3}”。写检验方案的格式是：取少量待测液于试管中，滴加某试剂，若观察到某现象，则说明某结论。回答“目的”或“原因”时说清“这样做是为了……，否则会……”，例如“趁热过滤是为了防止苯甲酸冷却结晶析出而损失”。设计定量方案要写清仪器、试剂、操作步骤、需要测量的数据和数据处理方法，例如2025年广东卷第17题(5) 用 pH 计测 0.100\molL{} NaClO 溶液的 pH，先求 \ce{ClO-} 的水解常数，再由 $K_\text{a}=\dfrac{K_\text{w}}{K_\text{h}}$ 求 HClO 的 $K_\text{a}$。

\begin{zhenti}{改编自2025年广东卷第17题(1)(2)(4)}
(1) 向盛有 2 mL 甲苯的试管中加入几滴酸性 \ce{KMnO4} 溶液，振荡，观察到体系颜色\tk。(2) 某苯甲酸粗品含少量泥沙和氯化钠，用重结晶法提纯，需要的操作及顺序为：加热溶解、\tk（从“a．冷却结晶，b．常温过滤，c．蒸馏，d．趁热过滤”中选择并排序）。(4) 甲同学根据“卤原子电负性越大，羧酸 $K_\text{a}$ 越大”推测邻氟苯甲酸（Ⅰ）$>$邻溴苯甲酸（Ⅱ）$>$苯甲酸（Ⅲ）的酸性，并测得常温下三种酸的饱和溶液 pH 大小顺序为Ⅲ$>$Ⅱ$>$Ⅰ，据此推断假设成立。乙同学认为该推断依据不足，因为\tk；乙同学用 pH 滴定法测得三种酸的 $K_\text{a}$ 顺序为Ⅱ$>$Ⅰ$>$Ⅲ。
\tcblower
\daan (1) 由紫色变为无色（甲苯被氧化为苯甲酸）。(2) dab（趁热过滤除泥沙，冷却结晶，常温过滤得晶体，NaCl 留在母液中）。(4) 常温下三种酸的饱和溶液浓度不同，不能直接用 pH 比较 $K_\text{a}$。\\
\dianbo 用 pH 比较酸性强弱，必须在浓度相同的前提下，这正是控制变量思想的应用。乙同学的实验结果也提醒我们，由已知规律外推得到的假设，必须用实验检验（邻位取代基的影响比较复杂）。
\end{zhenti}

% ---------------- 章末 ----------------
\zhangmo{本章小结}
\subsection*{一、知识总结}
\begin{center}
\small
\begin{tabularx}{\linewidth}{p{2.2cm}X}
\toprule
仪器操作 & 读数（按量筒分度值；滴定管 0.01 mL、托盘天平通常 0.1 g）；容量瓶、滴定管的使用；气密性检查；pH 试纸；洗涤与检验洗净；配制溶液 \\
分离检验 & 过滤（一贴二低三靠）、蒸发、蒸馏（水银球位置、沸石、冷凝水下进上出）、萃取分液（下层下口放、上层上口倒）、结晶与重结晶；离子检验关注干扰与顺序 \\
气体制备 & 发生 $\to$ 除杂 $\to$ 干燥 $\to$ 收集或性质实验 $\to$ 尾气处理；先除杂后干燥，洗气长进短出 \\
定量实验 & 滴定（关系式、电子守恒，误差归结到 $V(\text{标准})$）；量气（冷却、调平、平视）；重量法（恒重） \\
探究设计 & 控制变量与对照、空白实验；假设与验证；规范表述 \\
\bottomrule
\end{tabularx}
\end{center}
\subsection*{二、方法提升}
\begin{enumerate}
\item 判断装置能否达到目的：逐一检查原理是否正确、试剂是否合适、装置是否合理（加热、导管长短、防倒吸）、操作顺序是否正确。
\item 定量计算三步：写反应找关系式，代入数据（注意取样比例，如 25.00 mL/100.00 mL），计算结果（注意有效数字）。
\item 设计探究方案：先确定要比较的变量，其余条件保持相同；要证明“有剩余”，检验不足量的那种反应物。
\end{enumerate}
\subsection*{三、易错点归纳}
\begin{enumerate}
\item 用润湿的 pH 试纸测 pH；用 pH 试纸测氯水的 pH。
\item 检验 \ce{SO4^2-} 时用稀硝酸酸化，或先加 \ce{BaCl2}（\ce{SO3^2-}、\ce{Ag+} 会干扰）。
\item 气体实验中先干燥后除杂（除杂会重新带入水蒸气）。
\item 滴定管读数俯视、仰视的影响与量筒相反；滴定计算忘记换算取样比例，氧化还原滴定找错计量关系。
\item 量气时未冷却、未调平液面就读数。
\item 描述现象只写结论（如“生成 \ce{Fe(OH)3}”）而不写现象（“生成红褐色沉淀”）。
\end{enumerate}

\zhangmo{分层练习}
\zu{A组\quad 基础巩固}
\begin{xiti}[start=1]
\item 下列实验操作正确的是\xz
\xxa{用 10 mL 量筒量取 8.58 mL 稀盐酸}{用碱式滴定管盛装 \ce{KMnO4} 标准溶液}{蒸馏时，温度计水银球置于蒸馏烧瓶的支管口处}{分液时，上层液体也从下口放出}

\item 检验某溶液中是否含有 \ce{Fe^2+}（溶液中同时含有 \ce{Fe^3+}），最合适的试剂是\xz
\xx{KSCN 溶液}{氯水和 KSCN 溶液}{\ce{K3[Fe(CN)6]} 溶液}{NaOH 溶液}

\item \laiyuan{2024年广东卷第6题}\ 提纯 2.0 g 苯甲酸粗品（含少量 NaCl 和泥沙）的过程如下：苯甲酸粗品 $\xrightarrow{\text{加水加热}}$ 浊液 $\xrightarrow{\text{趁热过滤}}$ 滤液 $\xrightarrow{\text{X}}$ $\xrightarrow{\text{过滤洗涤}}$ 苯甲酸。其中，操作 X 为\xz
\xx{加热蒸馏}{加水稀释}{冷却结晶}{萃取分液}

\item 选择合适的分离方法（填序号：a．过滤\quad b．蒸发结晶\quad c．冷却结晶\quad d．萃取分液\quad e．蒸馏）：(1) 从食盐水中获得 NaCl\tk；(2) 从溴水中提取溴\tk；(3) 分离乙醇与水\tk；(4) 除去 \ce{KNO3} 中少量的 NaCl\tk。

\item 下列试剂的保存方法错误的是\xz
\xxa{浓硝酸保存在棕色试剂瓶中，置于阴凉处}{NaOH 溶液保存在带磨口玻璃塞的试剂瓶中}{少量金属钠保存在煤油中}{氢氟酸保存在塑料瓶中}

\item \laiyuan{原创}\ 实验室用 \ce{NH4Cl} 与 \ce{Ca(OH)2} 制取并收集干燥的 \ce{NH3}，下列说法错误的是\xz
\xxa{发生装置为固体加热型，试管口应略向下倾斜}{生成的 \ce{NH3} 可用浓硫酸干燥}{用向下排空气法收集 \ce{NH3}}{尾气用水吸收时，导管口接倒扣的漏斗以防倒吸}
\end{xiti}

\zu{B组\quad 高考典型}
\begin{xiti}
\item \laiyuan{改编自2025年广东卷第15题}\ 按图~\ref{fig:12-z25-15} 组装装置并进行实验：将铜丝插入溶液中，当 c 中红色褪去时，将铜丝拔离液面。下列叙述错误的是\xz
\tu{\begin{tikzpicture}[scale=0.85]
  \begin{scope}\clip (-0.28,2.3)--(-0.28,0.28) arc[start angle=180,end angle=360,radius=0.28]--(0.28,2.3)--cycle;
    \fill[c12liq] (-0.5,-0.1) rectangle (0.5,1.0);\end{scope}
  \draw[c12g] (-0.28,2.3)--(-0.28,0.28) arc[start angle=180,end angle=360,radius=0.28]--(0.28,2.3);
  \draw[c12stop] (-0.32,2.25) rectangle (0.32,2.45);
  \draw[c12wire] (-0.1,3.2)--(-0.1,1.25)--(-0.03,1.18)--(-0.17,1.1)--(-0.03,1.02)--(-0.17,0.94)--(-0.1,0.88);
  \draw[c12flow,<->] (-0.4,2.7)--(-0.4,3.2);
  \node[c12lab,left] at (-0.42,2.95) {可上下抽动\\的铜丝};
  \node[c12lab] at (0,-0.3) {a};
  \draw[c12ptr] (0.2,0.6)--(0.75,0.4) node[right,c12lab]{稀硝酸};
  \draw[c12t] (0.12,2.3)--(0.12,2.9)--(2.07,2.9)--(2.07,1.0);
  \draw[c12flow] (0.5,3.07)--(1.6,3.07);
  \pic at (2.2,0) {c12rbf};
  \node[c12lab] at (2.2,-0.3) {b};
  \draw[c12t] (2.33,1.6)--(2.33,2.9)--(4.2,2.9)--(4.2,0.35);
  \draw[c12flow] (2.7,3.07)--(3.8,3.07);
  \begin{scope}[xshift=4.2cm]
    \begin{scope}\clip (-0.28,2.3)--(-0.28,0.28) arc[start angle=180,end angle=360,radius=0.28]--(0.28,2.3)--cycle;
      \fill[red!12] (-0.5,-0.1) rectangle (0.5,1.2);\end{scope}
    \draw[c12g] (-0.28,2.3)--(-0.28,0.28) arc[start angle=180,end angle=360,radius=0.28]--(0.28,2.3);
    \draw[c12stop] (-0.32,2.25) rectangle (0.32,2.45);
    \node[c12lab] at (0,-0.3) {c};
    \draw[c12ptr] (0.2,0.7)--(0.7,0.5) node[right,c12lab]{含酚酞的\\\ce{NaHCO3} 溶液};
  \end{scope}
  \draw[c12t] (4.34,1.7)--(4.34,2.4)--(6.77,2.4);
  \draw[c12flow] (4.8,2.57)--(6.2,2.57);
  \pic at (6.9,0) {c12wash};
  \draw[c12t] (7.03,2.4)--(7.55,2.4);
  \node[c12lab] at (6.9,-0.3) {碱性 \ce{KMnO4} 溶液\\处理残余氮氧化物};
\end{tikzpicture}}{硝酸与铜的反应及气体性质实验\label{fig:12-z25-15}}
\xxa{a 中有化合反应发生，并有颜色变化}{b 中气体变红棕色时，所含氮氧化物至少有两种}{c 中溶液红色刚好褪去时，\ce{HCO3-} 恰好完全反应}{若将 a 中稀硝酸换为浓硫酸并加热，则 c 中溶液颜色会褪去}

\item \laiyuan{原创}\ 用如下方法测定某漂白粉中“有效氯”（以 \ce{Cl2} 计的质量分数）：称取 1.000 g 样品，加水研磨并充分浸取，过滤、洗涤，合并滤液和洗涤液后定容至 250.0 mL；取 25.00 mL，加入过量 KI 溶液和稀硫酸（\ce{ClO- + 2I- + 2H+ = Cl- + I2 + H2O}），用 0.1000\molL{} \ce{Na2S2O3} 标准溶液滴定生成的 \ce{I2}，近终点溶液呈浅黄色时加入淀粉溶液，消耗标准液共 9.00 mL。(1) 终点现象为\tk。(2) 求有效氯的质量分数。(3) 若滴定时间过长，空气中的 \ce{O2} 氧化了部分 \ce{I-}，结果偏\tk。

\item \laiyuan{改编自2026年广东卷第17题(1)～(3)}\ 某兴趣小组对铜的配合物进行探究实验。(1) 氨水的制备：从下列装置中选择需要的装置\tk。A．加热 NaOH 溶液与 \ce{NH4Cl} 固体的发生装置；B．烧杯中盛水、导管口接倒扣漏斗的吸收装置；C．盛 \ce{P2O5} 的干燥管；D．分液漏斗中盛甲苯与 \ce{NH4Cl} 水溶液的萃取装置。(2) 制备铜氨溶液，以其为电解液在铁件表面镀铜，铁件连接直流电源\tk{}极。(3) 铜氨配合物的制备与表征：向 \ce{CuSO4} 溶液中滴加氨水至过量，得深蓝色溶液，加入乙醇，析出固体，过滤，\tk，晾干得产物 Q。用 X 射线衍射仪测试 Q，观察到\tk，说明 Q 为晶体。

\item \laiyuan{原创，与第1章联系}\ 用 50 mL 0.50\molL{} 盐酸与 50 mL 0.55\molL{} NaOH 溶液测定中和反应的反应热，三次实验的温度差分别为 3.30 ℃、3.40 ℃、3.35 ℃（溶液密度按 1 g·mL$^{-1}$、比热容按 4.18 J·g$^{-1}$·℃$^{-1}$ 计算）。(1) 求测得的中和反应反应热 $\Delta H$（保留一位小数）。(2) 与理论值 $-57.3$\kjmol{} 相比，指出两条可能的原因。

\item \laiyuan{改编自2024年广东卷第17题(4)}\ 硫酸钙可用纯碱溶液处理：\ce{CaSO4(s) + CO3^2-(aq) <=> CaCO3(s) + SO4^2-(aq)}（反应Ⅰ）。兴趣小组将一定量硫酸钙加入 \ce{Na2CO3} 溶液中，搅拌并用 pH 计测体系 pH，一段时间后 pH 不再变化。甲同学认为反应Ⅰ已达到平衡；乙同学认为证据不足，提出假设1：硫酸钙固体已完全消耗；假设2：硫酸钙固体有剩余，但被碳酸钙沉淀包裹。乙同学将反应混合物固液分离，用少量蒸馏水快速洗去沉淀附着的母液（忽略洗涤时固体的溶解损失），取少量沉淀置于试管中，滴加\tk，观察到\tk，沉淀完全溶解；继续向试管中滴加\tk，无白色沉淀生成。由此得出假设1成立。(1) 补全上述空白。(2) 为什么“pH 不变”不足以说明反应Ⅰ达到平衡？(3) 为什么要先使沉淀完全溶解再检验？
\end{xiti}

\zu{C组\quad 冲刺高分}
\begin{xiti}
\item \laiyuan{原创}\ 某小组探究 \ce{FeCl3} 与 \ce{Na2SO3} 溶液混合后的变化：体系先由黄色变为红褐色，无明显气体和沉淀产生；放置后逐渐变为浅绿色。已知除水解、氧化还原外，\ce{Fe^3+} 也可能与 \ce{SO3^2-} 形成有色配合物，试剂中原本不含 \ce{Fe^2+} 和 \ce{SO4^2-}。
(1) 甲认为“出现红褐色就证明生成了 \ce{Fe(OH)3} 胶体”。这个判断是否充分？可用什么现象初步区分胶体与一般溶液？
(2) 若发生 \ce{Fe^3+} 氧化 \ce{SO3^2-} 的反应，写出离子方程式。
(3) 设计实验证明最终体系中存在 \ce{Fe^2+} 和 \ce{SO4^2-}（分别取样，写出操作、现象和结论）。
(4) 仅凭“先红褐色、后浅绿色”，能否断言“水解比氧化还原快，且前者的反应限度较小”？说明理由，并写出进一步研究速率需要测量的量。

\item \laiyuan{原创}\ 铜与浓硝酸反应后的溶液呈绿色，而铜与稀硝酸反应后的溶液呈蓝色。甲同学认为是 \ce{Cu(NO3)2} 浓度较大所致；乙同学认为是溶液中溶解了 \ce{NO2}（\ce{NO2} 溶于 \ce{Cu(NO3)2} 溶液呈黄色，与蓝色混合呈绿色）。(1) 设计实验验证甲的假设（写出操作和预期现象）。(2) 设计实验验证乙的假设。(3) 若实验证实乙的假设正确，如何使绿色溶液转变为蓝色？
\end{xiti}

\zhangmo{练习解析}
\begin{jiexilist}
\item \ans{C}A：量筒精度为 0.1 mL，不能量取 8.58 mL；B：\ce{KMnO4} 有强氧化性，会腐蚀碱式滴定管的橡胶管，应用酸式滴定管；C 正确；D：上层液体应从上口倒出。
\item \ans{C}\ce{Fe^3+} 存在时，加 KSCN 溶液本身就变红，无法判断 \ce{Fe^2+}；NaOH 溶液生成的红褐色沉淀会掩盖白色沉淀的颜色变化；\ce{K3[Fe(CN)6]} 与 \ce{Fe^2+} 生成蓝色沉淀，不受 \ce{Fe^3+} 干扰。
\item \ans{C}苯甲酸的溶解度随温度变化大，NaCl 变化小。趁热过滤除去泥沙后，冷却结晶使苯甲酸析出，NaCl 留在溶液中。
\item (1) b；(2) d（用 \ce{CCl4} 或苯萃取）；(3) e（普通蒸馏可富集乙醇；要制得无水乙醇，通常先用过量 CaO 除水再蒸馏）；(4) c（\ce{KNO3} 溶解度随温度变化大，NaCl 变化小，冷却时主要析出 \ce{KNO3}）。
\item \ans{B}NaOH 与玻璃中的 \ce{SiO2} 反应生成有黏性的 \ce{Na2SiO3}，会使瓶塞与瓶口粘在一起，应用橡胶塞。
\item \ans{B}\ce{NH3} 是碱性气体，会与浓硫酸反应（\ce{2NH3 + H2SO4 = (NH4)2SO4}），应用碱石灰干燥。A：固体加热时管口略向下倾斜，防止冷凝水倒流使试管炸裂；C：\ce{NH3} 密度比空气小；D：\ce{NH3} 极易溶于水，需防倒吸。
\item \ans{C}A：Cu 与稀硝酸反应生成 NO，NO 与 a 中液面上方空气中的 \ce{O2} 化合生成红棕色 \ce{NO2}，溶液变蓝，正确。B：b 中出现红棕色说明有 \ce{NO2}，而 \ce{NO2} 存在 \ce{2NO2 <=> N2O4}，所以氮氧化物至少有两种，正确。C：酚酞的变色范围为 pH 8.2～10.0，红色褪去时溶液可能仍呈弱碱性，不能说明 \ce{HCO3-} 恰好完全反应，错误。D：浓硫酸与铜加热生成 \ce{SO2}，\ce{SO2} 通入 \ce{NaHCO3} 溶液生成 \ce{NaHSO3}，溶液呈酸性，红色褪去，正确。
\item (1) 滴入最后半滴 \ce{Na2S2O3} 溶液时，溶液蓝色褪去，且半分钟内不恢复。(2) 有效氯的定义是样品与盐酸作用放出的 \ce{Cl2} 的质量分数，由 \ce{ClO- + Cl- + 2H+ = Cl2 ^ + H2O} 知 1 mol \ce{ClO-} 相当于 1 mol \ce{Cl2}，关系式为 \ce{Cl2} $\sim$ \ce{ClO-} $\sim$ \ce{I2} $\sim$ 2\ce{S2O3^2-}。$n(\ce{S2O3^2-})=9.000\times10^{-4}$ mol，25.00 mL 中 $n(\ce{Cl2})=4.500\times10^{-4}$ mol，250.0 mL 中共 $4.500\times10^{-3}$ mol，$m(\ce{Cl2})=4.500\times10^{-3}\times71=0.3195$ g，有效氯约为 32.0\%。(3) 高。\ce{O2} 氧化 \ce{I-} 生成额外的 \ce{I2}，消耗的 \ce{Na2S2O3} 偏多。
\item (1) AB。用 NaOH 与 \ce{NH4Cl} 加热制 \ce{NH3}（A），用倒扣漏斗将 \ce{NH3} 溶于水，防倒吸（B）；\ce{P2O5} 是酸性干燥剂，会与 \ce{NH3} 反应；D 是萃取装置，与制氨水无关。(2) 负（镀件作阴极，\ce{Cu} 在铁件表面析出）。(3) 洗涤；明锐的衍射峰（或分立的斑点）。
\item (1) 平均温差 $\Delta t=\dfrac{3.30+3.40+3.35}{3}=3.35$ ℃；$Q=100\times4.18\times3.35\approx1400$ J$=1.40$ kJ；生成水 0.025 mol，$\Delta H=-\dfrac{1.40}{0.025}\approx-56.0$\kjmol。(2) 测得放出的热量偏少，可能原因：装置隔热效果不好，有热量散失；NaOH 溶液分多次倒入；测完酸的温度后未冲洗温度计就测碱的温度；溶液比热容、密度取近似值带来误差等（任答两条）。
\item (1) 过量稀盐酸；有气体产生；\ce{BaCl2} 溶液。(2) pH 不变只说明 \ce{CO3^2-} 浓度不再变化，可能是正、逆反应速率相等（平衡），也可能是反应物之一已经耗尽、反应停止了，两种情况 pH 都不变。(3) 若沉淀中有被 \ce{CaCO3} 包裹的 \ce{CaSO4}，只有使沉淀完全溶解，\ce{SO4^2-} 才能全部进入溶液被检验到；先加过量盐酸还可除去 \ce{CO3^2-}，避免生成 \ce{BaCO3} 沉淀干扰。洗去附着母液可排除原溶液中 \ce{SO4^2-} 的干扰；加 \ce{BaCl2} 无白色沉淀，在本题模型和检出范围内支持固体已无 \ce{CaSO4}，即假设1。
\item (1) 不充分。红褐色也可能来自有色配合物，颜色不能唯一确定微粒；可用一束光照射，观察有无明显丁达尔效应作初步区分，仍需结合其他证据。(2) \ce{2Fe^3+ + SO3^2- + H2O = 2Fe^2+ + SO4^2- + 2H+}，两边电荷均为 $+4$。(3) 取少量最终溶液，加 \ce{K3[Fe(CN)6]} 溶液产生蓝色沉淀，说明有 \ce{Fe^2+}；另取一份，先加足量稀盐酸，待气体逸尽，再加 \ce{BaCl2}，出现白色沉淀，说明有 \ce{SO4^2-}。不用硝酸酸化，以免把剩余亚硫酸根氧化为硫酸根。(4) 不能。红褐色微粒尚未确定，颜色先后只说明可见变化的先后，不能直接给出两个反应的速率或平衡转化程度。应在相同条件下测定特定物质浓度随时间的变化，例如 $c(\ce{Fe^2+})$ 随时间的变化；判断限度还要在确认平衡后测平衡组成，不能把“变色结束”直接当作平衡。
\item (1) 配制不含氮氧化物、而 \ce{Cu^2+} 浓度及硝酸浓度均与待测液相近的对照液，保持相同温度，用相同容器、相同液层厚度比较颜色：若对照液仍蓝，说明仅用铜离子浓度不能解释绿色；若也绿，则支持浓度影响的假设，但还应排除其他因素。(2) 对少量绿色溶液通入 \ce{N2} 并温和加热，同时给等量蓝色对照液作相同处理；逸出的气体另行吸收。若绿色样品转蓝而对照液无相同变化，支持乙的假设；也可在蓝色对照液中溶入少量 \ce{NO2}，观察是否重现绿色。(3) 按(2)驱除 \ce{NO2}，或通入 \ce{O2} 使其转化，\ce{4NO2 + O2 + 2H2O = 4HNO3}。单纯加水稀释同时改变铜离子浓度及氮氧化物状态，不能单独作为区分两个假设的证据。
\end{jiexilist}
